\documentclass[12pt, a4paper]{report}

%% ==================== PACKAGES ====================
\usepackage{fontspec}
\usepackage{polyglossia}
\usepackage{geometry}
\usepackage{xcolor}
\usepackage{colortbl}
\usepackage{graphicx}
\usepackage{overpic}
\usepackage{booktabs}
\usepackage{array}
\usepackage{multirow}
\usepackage{longtable}
\usepackage{setspace}
\usepackage{fancyhdr}
\usepackage{eso-pic}
\usepackage{amsmath}
\usepackage{enumitem}
\usepackage{tikz}
\usepackage{ifthen}
\usepackage{tocloft}
\usepackage{caption}
\usepackage{hyperref}
\usepackage{silence}
\WarningFilter{latex}{File}

%% ==================== LANGUAGE ====================
\setmainlanguage[numerals=maghrib]{arabic}
\setotherlanguage{english}
\newcommand{\lr}[1]{\textenglish{#1}}

%% ==================== FONTS ====================
\newfontfamily\arabicfont[Script=Arabic, Scale=1.1]{Amiri}
\newfontfamily\arabicfontsf[Script=Arabic]{Amiri}
\newfontfamily\arabicfonttt[Script=Arabic]{Amiri}
\setmainfont{Amiri}[Script=Arabic]

%% ==================== PAGE GEOMETRY ====================
\geometry{
  a4paper,
  top=3cm,
  bottom=2.5cm,
  left=3cm,
  right=3.5cm,
  headheight=1.5cm
}

%% ==================== LINE SPACING ====================
\setstretch{1.5}

%% ==================== COLORS ====================
\definecolor{primaryblue}{RGB}{0, 51, 102}
\definecolor{sectiongray}{RGB}{64, 64, 64}
\definecolor{tableheader}{RGB}{220, 230, 242}
\newif\ifdrawborder
\drawbordertrue
% If logo files don't exist, use placeholder boxes
\newcommand{\logoplaceholder}[2]{%
  \fbox{\parbox[c][#1][c]{#1}{\centering\tiny #2}}%
}

%% ==================== HEADER/FOOTER ====================
\pagestyle{fancy}
\fancyhf{}
\fancyhead[R]{\small\color{sectiongray}\RL{تأثير التسويق الإلكتروني على سياسة المبيعات}}
\fancyhead[L]{\small\color{sectiongray}\LR{\thepage}}
\fancyfoot[C]{\small\color{sectiongray}\RL{جامعة البصرة للنفط والغاز --- كلية الإدارة الصناعية}}
\renewcommand{\headrulewidth}{0.5pt}
\renewcommand{\footrulewidth}{0.3pt}

\AddToShipoutPictureFG{%
  \ifdrawborder
  \begin{tikzpicture}[remember picture, overlay]
    \draw[line width=1.5pt, color=primaryblue]
      ([xshift=0.7cm, yshift=-0.7cm]current page.north west)
      rectangle
      ([xshift=-0.7cm, yshift=0.7cm]current page.south east);
    \draw[line width=0.5pt, color=primaryblue!60]
      ([xshift=0.9cm, yshift=-0.9cm]current page.north west)
      rectangle
      ([xshift=-0.9cm, yshift=0.9cm]current page.south east);
  \end{tikzpicture}%
  \fi
}

%% ==================== CHAPTER/SECTION TITLES ====================
\makeatletter
\def\@makechapterhead#1{%
  \vspace*{50\p@}%
  {\parindent \z@ \centering \normalfont
    \if@RTL\beginR\fi
    \ifnum \c@secnumdepth >\m@ne
      {\huge\bfseries\color{primaryblue} \@chapapp\ \thechapter\par\nobreak}
      \vskip 20\p@
    \fi
    {\Huge\bfseries\color{primaryblue} #1\par\nobreak}
    \if@RTL\endR\fi
    \vskip 40\p@
  }}

\def\@makeschapterhead#1{%
  \vspace*{50\p@}%
  {\parindent \z@ \centering \normalfont
    \if@RTL\beginR\fi
    {\Huge\bfseries\color{primaryblue} #1\par\nobreak}
    \if@RTL\endR\fi
    \vskip 40\p@
  }}

\renewcommand\section{\@startsection{section}{1}{\z@}%
  {-3.5ex \@plus -1ex \@minus -.2ex}%
  {2.3ex \@plus .2ex}%
  {\normalfont\Large\bfseries\color{primaryblue}}}

\renewcommand\subsection{\@startsection{subsection}{2}{\z@}%
  {-3.25ex\@plus -1ex \@minus -.2ex}%
  {1.5ex \@plus .2ex}%
  {\normalfont\large\bfseries\color{sectiongray}}}

\renewcommand\subsubsection{\@startsection{subsubsection}{3}{\z@}%
  {-3.25ex\@plus -1ex \@minus -.2ex}%
  {1.5ex \@plus .2ex}%
  {\normalfont\normalsize\bfseries\color{sectiongray}}}
\makeatother

%% ==================== TOC DEPTH ====================
\setcounter{tocdepth}{3}
\setcounter{secnumdepth}{3}

%% ==================== HYPERREF ====================
\hypersetup{
  colorlinks=true,
  linkcolor=primaryblue,
  urlcolor=primaryblue,
  citecolor=primaryblue,
  pdftitle={تأثير التسويق الإلكتروني على سياسة المبيعات},
  pdfauthor={محمد مازن مكي}
}

%% ==================== MISC ====================
\setlength{\parindent}{1.5em}
\setlength{\parskip}{6pt}
\let\cleardoublepage\clearpage

%%============================================================
\begin{document}
%%============================================================

\pagenumbering{roman}
\pagestyle{empty}

%%------------------------------------------------------------
%% TITLE PAGE
%%------------------------------------------------------------
\begin{titlepage}

%% --- Corner Logos ---
\begin{tikzpicture}[remember picture, overlay]
  %% Top-right: University logo
  \node[anchor=north east, inner sep=1cm]
    at (current page.north east)
    {\IfFileExists{cool.jpg}%
      {\includegraphics[width=5cm]{cool.jpg}}%
      {\logoplaceholder{5cm}{شعار الجامعة}}};
  %% Top-left: College logo
  \node[anchor=north west, inner sep=1.5cm]
    at (current page.north west)
    {\IfFileExists{cool2.jpg}%
      {\includegraphics[width=4cm]{cool2.jpg}}%
      {\logoplaceholder{2.5cm}{شعار الكلية}}};
\end{tikzpicture}

\begin{center}

\vspace*{0.5cm}
{\large\bfseries وزارة التعليم العالي والبحث العلمي}\\[0.2cm]
{\large\bfseries جامعة البصرة للنفط والغاز}\\[0.2cm]
{\large\bfseries كلية الإدارة الصناعية للنفط والغاز}\\[0.2cm]
{\large\bfseries قسم إدارة وتسويق النفط والغاز}

\vspace{0.8cm}
\noindent\rule{\linewidth}{1.5pt}
\vspace{0.4cm}

{\huge\bfseries\color{primaryblue}
تأثير التسويق الإلكتروني\\[0.4cm]
على سياسة المبيعات\\[0.4cm]
في المؤسسات التجارية}

\vspace{0.6cm}
{\large دراسة حالة: شركة توزيع المنتجات النفطية --- البصرة}

\vspace{0.8cm}
\noindent\rule{\linewidth}{0.7pt}
\vspace{0.4cm}

{\large
\begin{tabular}{rr}
\textbf{بحث مقدم إلى:} & مجلس كلية الإدارة الصناعية للنفط والغاز \\[0.2cm]
\textbf{وهو جزء من:} & متطلبات نيل شهادة البكالوريوس \\[0.2cm]
                      & في علوم إدارة وتسويق النفط والغاز \\
\end{tabular}
}

\vspace{0.2cm}

\begin{tabular}{rr}
\textbf{\large إعداد الطالب:}  & \textbf{\large محمد مازن مكي} \\[0.5cm]
\textbf{\large بإشراف:}        & \textbf{\large المدرس المساعد ندى عدنان جواد} \\
\end{tabular}
\vspace{0.2cm}
\vfill
{\large\bfseries 2026م \quad---\quad 1447هـ}

\end{center}
\end{titlepage}

%%------------------------------------------------------------
%% QURANIC VERSE
%%------------------------------------------------------------
\newpage\thispagestyle{empty}
\vspace*{4cm}
\begin{center}
{\Large\bfseries بسم الله الرحمن الرحيم}\\[1cm]
{\large
«هُوَ الَّذِي جَعَلَ الشَّمْسَ ضِيَاءً وَالْقَمَرَ نُورًا وَقَدَّرَهُ مَنَازِلَ لِتَعْلَمُوا عَدَدَ السِّنِينَ وَالْحِسَابَ ۚ مَا خَلَقَ اللَّهُ ذَٰلِكَ إِلَّا بِالْحَقِّ ۚ يُفَصِّلُ الْآيَاتِ لِقَوْمٍ يَعْلَمُونَ»\\[0.5cm]
(سورة يونس: الآية 5)
}\\[0.8cm]
{\large\bfseries صدق الله العظيم}
\end{center}
\newpage
%%------------------------------------------------------------
%% SUPERVISOR CERTIFICATE
%%------------------------------------------------------------
\newpage\thispagestyle{empty}
\chapter*{إقرار المشرف}
\addcontentsline{toc}{chapter}{إقرار المشرف}

\noindent بسم الله الرحمن الرحيم

\vspace{0.5cm}
\noindent أشهد أن إعداد بحث التخرج هذا والموسوم:

\begin{center}
\fbox{\parbox{0.85\textwidth}{\centering\vspace{6pt}
\textbf{تأثير التسويق الإلكتروني على سياسة المبيعات في المؤسسات التجارية}\\[4pt]
\textbf{دراسة حالة: شركة توزيع المنتجات النفطية --- البصرة}
\vspace{6pt}}}
\end{center}

\vspace{0.5cm}
\noindent للطالب: \textbf{محمد مازن مكي}، قد جرى تحت إشرافي في كلية الإدارة الصناعية للنفط والغاز / قسم إدارة وتسويق النفط والغاز (البكالوريوس)، للعام الدراسي 2025--2026، وهو جزء من متطلبات نيل شهادة البكالوريوس في تخصص إدارة وتسويق النفط والغاز.

\vspace{0.8cm}

\begin{tabular}{r l}
\textbf{توقيع المشرف:} & \underline{\hspace{6cm}} \\[0.8cm]
\textbf{اسم المشرف:} & المدرس المساعد ندى عدنان جواد \\[0.8cm]
\textbf{التاريخ:} & \underline{\hspace{6cm}} \\
\end{tabular}

\vspace{1cm}

\noindent\textbf{توصية رئيس القسم العلمي (إدارة وتسويق النفط والغاز):}

\vspace{0.3cm}
\noindent بناءً على توصية المشرف العلمي، أرشِّح هذا البحث للمناقشة.

\vspace{1cm}

\begin{tabular}{r l}
\textbf{اسم رئيس القسم:} & الأستاذ المساعد الدكتور أحمد علي أحمد \\[0.5cm]
\textbf{التوقيع:} & \underline{\hspace{6cm}} \\[0.5cm]
\textbf{التاريخ:} & \quad /\quad /\quad \\
\end{tabular}

%%------------------------------------------------------------
%% DEDICATION
%%------------------------------------------------------------
\newpage\thispagestyle{empty}
\chapter*{الإهداء}
\addcontentsline{toc}{chapter}{الإهداء}

بعد مسيرةٍ حافلة بالتحديات والتضحيات، ورغم ما مررنا به من صعوبات وآلام، ها نحن نبلغ ختام رحلتنا الدراسية.

أتقدم بإهداء هذا المشروع إلى \textbf{أسرتي الكريمة}، التي كانت العون والسند طوال هذا المشوار، وإلى أخواتي العزيزات، وإلى أصدقائي الذين كانوا نعم الرفقة والدعم، وكانوا بمثابة إخوة في كل الظروف.

كما أُعرب عن بالغ امتناني لكل من قدَّم لي الدعم والمساندة، وأسهم في استمراري وتقدمي.

وأخص بالشكر والتقدير الكادر التدريسي الكريم، لما قدمه من علم وتوجيه وإلهام.

وختاماً، أُهدي ثمرة هذا الجهد المتواضع إلى كل من كان له أثرٌ في وصولي إلى هذه المرحلة، مع أصدق الأمنيات بالتوفيق للجميع.

\vspace{1cm}
\begin{flushright}
\textbf{الباحث: محمد مازن مكي}
\end{flushright}

%%------------------------------------------------------------
%% ACKNOWLEDGMENTS
%%------------------------------------------------------------
\newpage\thispagestyle{empty}
\chapter*{شكر وتقدير}
\addcontentsline{toc}{chapter}{شكر وتقدير}

الحمد لله أولاً وآخراً، الذي بنعمته تتم الصالحات، وبتوفيقه يتحقق النجاح.

أتقدم بجزيل الشكر والتقدير إلى كل من كان له دور في دعمي ومساندتي طوال مسيرتي العلمية، وإلى \textbf{والدتي الحبيبة}، التي كانت وما زالت رمزاً للعطاء غير المشروط، وملاذاً للأمان والحنان، أهديكِ يا أماه كل الامتنان والوفاء، فلكِ الفضل الأول بعد الله فيما وصلتُ إليه اليوم.

كما لا يفوتني أن أتوجه بخالص الشكر والتقدير إلى \textbf{والدي العزيز}، الذي قدَّم لي من الدعم والتشجيع ما لا يُقدَّر بثمن، وكان دائماً السند والمعين في كل مراحل حياتي.

إلى أساتذتي الكرام، الذين كان لهم الأثر الكبير في مسيرتي التعليمية، أخص بالشكر من وجَّهني، وصبر على تعليمي، وفتح لي آفاق العلم والمعرفة، فلهم مني كل الاحترام والتقدير.

وأتوجه بخالص الشكر إلى أستاذتي المشرفة \textbf{المدرس المساعد ندى عدنان جواد}، على توجيهاتها السديدة وصبرها في الإشراف على هذا البحث.

إلى أصدقائي الأعزاء، الذين شاركوني الرحلة بكل ما فيها من تحديات وأوقات لا تُنسى، لكم مني وافر الامتنان والمحبة على وقوفكم إلى جانبي ودعمكم الدائم.

وإلى كل من قدَّم لي يد العون، أو كلمة طيبة، أو دعاءً صادقاً، أو وقف بجانبي في أوقات الشدة، أقول: شكراً لكم من القلب، فأنتم شركاء هذا الإنجاز.

\vspace{1cm}
\begin{flushright}
\textbf{الباحث: محمد مازن مكي}
\end{flushright}

%%------------------------------------------------------------
%% TABLE OF CONTENTS & LIST OF TABLES
%%------------------------------------------------------------
\newpage
\tableofcontents
\newpage
\listoftables

%%============================================================
%%  MAIN MATTER
%%============================================================
\newpage
\pagenumbering{arabic}
\pagestyle{fancy}
\drawbordertrue

%%============================================================
%%  CHAPTER ONE: RESEARCH METHODOLOGY
%%============================================================
\chapter{منهجية البحث والدراسات السابقة}

\section{المبحث الأول: منهجية البحث}

\subsection{مقدمة}

شهد العالم في السنوات الأخيرة تطوراً كبيراً في مجال تكنولوجيا المعلومات والاتصالات، مما أدى إلى ظهور أنماط جديدة من التعاملات التجارية عُرفت بالتجارة الإلكترونية. وقد أصبح التسويق الإلكتروني أحد أبرز المجالات التي استفادت من هذا التطور، إذ غيَّر بشكل جذري أساليب الترويج والتواصل مع العملاء. تسعى المؤسسات التجارية اليوم إلى مواكبة هذا التحول الرقمي من خلال تبنِّي استراتيجيات تسويق إلكترونية تُسهم في تعزيز قدرتها التنافسية وزيادة مبيعاتها في الأسواق المحلية والدولية.

وفي قطاع النفط والغاز تحديداً، باتت شركات التوزيع والتسويق النفطي أمام ضرورة ملحَّة لتبنِّي الأدوات الرقمية في إدارة علاقاتها مع العملاء وتعزيز قدرتها التنافسية. وفي ظل هذا التحول، أصبح من الضروري دراسة تأثير التسويق الإلكتروني على سياسات المبيعات داخل المؤسسات، ومدى إسهامه في تحسين الأداء البيعي وتطوير العلاقة مع العملاء.

أدى التطور السريع في استخدام الإنترنت ووسائل التواصل الاجتماعي إلى تحويل التسويق من مفهومه التقليدي إلى مفهوم أكثر ديناميكية يعتمد على التفاعل المباشر مع المستهلكين عبر المنصات الرقمية. وقد باتت المؤسسات التجارية تستخدم أدوات متنوعة مثل المواقع الإلكترونية، والإعلانات الرقمية، والتسويق عبر البريد الإلكتروني، والتسويق بالمحتوى للوصول إلى جمهورها المستهدف بطرق أكثر فاعلية وأقل تكلفة.

ومن هنا جاءت هذه الدراسة لتُحلِّل واقع تطبيق التسويق الإلكتروني في شركة توزيع المنتجات النفطية بالبصرة، وأثره على سياسة مبيعاتها.

\subsection{مشكلة البحث}

تتمثل مشكلة هذا البحث في التحول الرقمي المتسارع الذي يشهده العالم، والاعتماد المتزايد للمؤسسات التجارية على أدوات التسويق الإلكتروني بصورة كلية أو شبه كلية. وقد أثار هذا التحول التساؤلات الآتية:

\begin{enumerate}[label=\arabic*.]
  \item ما مدى قدرة التسويق الإلكتروني على التأثير في زيادة المبيعات داخل شركة توزيع المنتجات النفطية؟
  \item ما طبيعة الأثر الذي يتركه التسويق الإلكتروني على سياسات التسعير وآليات البيع؟
  \item كيف يؤثر التسويق الإلكتروني في سلوك المستهلكين، سواء أكان إيجابياً أم سلبياً، ومدى إسهامه في تحسين أو إضعاف مستوى رضاهم عن المنتجات والخدمات المقدمة؟
  \item ما دور أدوات التسويق الرقمي في تحقيق الميزة التنافسية في بيئة الأعمال النفطية؟
\end{enumerate}

\subsection{أهمية البحث}

تكمن أهمية هذا البحث في جوانب عدة، أبرزها:

\begin{itemize}
  \item \textbf{الأهمية العلمية:} يُعد هذا البحث من الدراسات النادرة التي تربط بين التسويق الإلكتروني وسياسة المبيعات في القطاع النفطي العراقي، مما يُثري الأدبيات العربية في هذا المجال. كما يُسهم في إثراء الأدبيات العربية من خلال الربط بين أدوات التسويق الإلكتروني وعناصر سياسة المبيعات بصورة علمية ومنهجية.
  \item \textbf{الأهمية التطبيقية:} يُساعد المؤسسات التجارية وموظفي المبيعات على اتخاذ قرارات دقيقة مبنية على التحليل الإحصائي للبيانات، مما يعزز من قدرتهم على تحسين الأداء البيعي وتطوير خططهم التسويقية.
  \item \textbf{الأهمية الاستراتيجية:} يُوفِّر نموذجاً تحليلياً يعتمد على بيانات دقيقة يتم جمعها وتحليلها لقياس تأثير التسويق الإلكتروني في سياسة المبيعات داخل شركة توزيع المنتجات النفطية.
\end{itemize}

\subsection{أهداف البحث}

يهدف هذا البحث إلى تحقيق الأهداف الآتية:

\begin{enumerate}[label=\arabic*.]
  \item قياس دور وأثر التسويق الإلكتروني في زيادة المبيعات داخل شركة توزيع المنتجات النفطية.
  \item تحليل تأثير التسويق الإلكتروني على سياسة المبيعات وطرق البيع في الأسواق.
  \item تشخيص المشكلات المحتملة التي قد تواجه الشركة عند استخدام أدوات التسويق الإلكتروني، من خلال تحليل البيانات والمقارنة بين النتائج.
  \item اقتراح حلول عملية تُساعد المؤسسات على اتخاذ قرارات مبنية على الأدلة والإحصاءات.
  \item تحديد مفهوم التسويق الإلكتروني وعناصره، وبيان مدى تأثيره في السلوك الشرائي ودرجة الولاء للعلامة التجارية.
\end{enumerate}

\subsection{فرضيات البحث}

انطلاقاً من مشكلة البحث وأهدافه، صِيغت الفرضيات الآتية:

\subsubsection{الفرضية الرئيسية}

\begin{itemize}
  \item[$H_0$:] \textbf{لا يوجد أثر ذو دلالة إحصائية} عند مستوى معنوية ($\alpha \leq 0.05$) للتسويق الإلكتروني على سياسة المبيعات في شركة توزيع المنتجات النفطية.
  \item[$H_1$:] \textbf{يوجد أثر ذو دلالة إحصائية} عند مستوى معنوية ($\alpha \leq 0.05$) للتسويق الإلكتروني على سياسة المبيعات في شركة توزيع المنتجات النفطية.
\end{itemize}

\subsubsection{الفرضيات الفرعية}

\begin{enumerate}[label=\arabic*.]
  \item
    \begin{itemize}
      \item[$H_{0_1}$:] لا يوجد أثر ذو دلالة إحصائية للتسويق عبر وسائل التواصل الاجتماعي على حجم المبيعات.
      \item[$H_{1_1}$:] يوجد أثر ذو دلالة إحصائية للتسويق عبر وسائل التواصل الاجتماعي على حجم المبيعات.
    \end{itemize}
  \item
    \begin{itemize}
      \item[$H_{0_2}$:] لا يوجد أثر ذو دلالة إحصائية للإعلانات الرقمية المدفوعة على سياسة التسعير.
      \item[$H_{1_2}$:] يوجد أثر ذو دلالة إحصائية للإعلانات الرقمية المدفوعة على سياسة التسعير.
    \end{itemize}
  \item
    \begin{itemize}
      \item[$H_{0_3}$:] لا يوجد أثر ذو دلالة إحصائية للتسويق عبر البريد الإلكتروني على رضا العملاء.
      \item[$H_{1_3}$:] يوجد أثر ذو دلالة إحصائية للتسويق عبر البريد الإلكتروني على رضا العملاء.
    \end{itemize}
  \item
    \begin{itemize}
      \item[$H_{0_4}$:] لا يوجد أثر ذو دلالة إحصائية لتحسين محركات البحث \lr{(SEO)} على قنوات التوزيع.
      \item[$H_{1_4}$:] يوجد أثر ذو دلالة إحصائية لتحسين محركات البحث \lr{(SEO)} على قنوات التوزيع.
    \end{itemize}
\end{enumerate}

\subsection{نموذج الدراسة}

يُوضح الجدول الآتي النموذج الهيكلي للعلاقة بين المتغيرات:

\begin{center}
\begin{tabular}{|c|c|c|}
\hline
\rowcolor{tableheader}
\textbf{المتغير المستقل (أبعاد التسويق الإلكتروني)} & \textbf{العلاقة} & \textbf{المتغير التابع} \\
\hline
التسويق عبر وسائل التواصل الاجتماعي & \multirow{4}{*}{$\longrightarrow$} & \multirow{4}{*}{سياسة المبيعات} \\
\cline{1-1}
الإعلانات الرقمية المدفوعة & & \\
\cline{1-1}
التسويق بالبريد الإلكتروني & & \\
\cline{1-1}
تحسين محركات البحث \lr{(SEO)} & & \\
\hline
\end{tabular}
\end{center}

\subsection{منهج الدراسة}

اعتمد هذا البحث على \textbf{المنهج الوصفي التحليلي}، لملاءمته لطبيعة موضوع الدراسة، إذ يُسهم في وصف الظاهرة كما هي في الواقع، ومن ثم تحليلها علمياً ومنطقياً للكشف عن العلاقة بين التسويق الإلكتروني وسياسة المبيعات في شركة توزيع المنتجات النفطية. وقد تم اختيار هذا المنهج لكونه لا يعتمد على التجربة المباشرة، بل يرتكز على جمع البيانات والمعلومات من مصادرها المختلفة، ومن ثم تحليلها وتفسيرها للوصول إلى نتائج دقيقة تدعم أهداف البحث.

\subsection{أداة الدراسة}

اعتمدت الدراسة على \textbf{الاستبانة} أداةً رئيسيةً لجمع البيانات الأولية، وقد صُمِّمت وفق مقياس ليكرت الخماسي \lr{(Likert Scale)}، وتضمنت المحاور الآتية:

\begin{enumerate}[label=\arabic*.]
  \item \textbf{المحور الأول:} البيانات الشخصية (الجنس، المؤهل العلمي، سنوات الخبرة، القسم الوظيفي).
  \item \textbf{المحور الثاني:} مستوى تطبيق أدوات التسويق الإلكتروني في الشركة.
  \item \textbf{المحور الثالث:} تأثير التسويق الإلكتروني على سياسة المبيعات.
  \item \textbf{المحور الرابع:} التحديات والمعوقات التي تواجه تطبيق التسويق الإلكتروني.
\end{enumerate}

\subsection{حدود وعينة البحث}

\subsubsection{أولاً: حدود البحث}

\begin{itemize}
  \item \textbf{الحدود الزمانية:} يقتصر البحث على الفترة الممتدة من (2020--2024)، وهي مرحلة شهدت تنامياً ملحوظاً في تأثير التسويق الإلكتروني على سياسة المبيعات في المؤسسات التجارية.
  \item \textbf{الحدود المكانية:} يقتصر البحث على شركة توزيع المنتجات النفطية في البصرة بوصفها دراسة حالة، لكونها إحدى الشركات النفطية العراقية التي تطبق استراتيجيات التسويق الصناعي.
  \item \textbf{الحدود الموضوعية:} التسويق الإلكتروني بوصفه متغيراً مستقلاً، وسياسة المبيعات بوصفها متغيراً تابعاً.
\end{itemize}

\subsubsection{ثانياً: مجتمع وعينة البحث}

يتمثل \textbf{مجتمع البحث} في جميع العاملين في شركة توزيع المنتجات النفطية من المدراء والموظفين، ولا سيما في الأقسام ذات العلاقة مثل إدارة التسويق والتخطيط والمبيعات.

أما \textbf{عينة البحث}، فتتمثل في عدد من المدراء والموظفين العاملين في الشركة، وتم اختيارها بالأسلوب العشوائي الطبقي، وبلغ حجمها (\textbf{60 مفردة}) وُزِّعت على النحو الآتي:

\begin{center}
\begin{tabular}{|r|c|c|}
\hline
\rowcolor{tableheader}
\textbf{القسم} & \textbf{العدد} & \textbf{النسبة \%} \\
\hline
قسم التسويق والتخطيط & 25 & 41.7\% \\
\hline
قسم المبيعات والتوزيع & 20 & 33.3\% \\
\hline
الإدارة العليا والوسطى & 15 & 25.0\% \\
\hline
\textbf{المجموع} & \textbf{60} & \textbf{100\%} \\
\hline
\end{tabular}
\end{center}

%%============================================================
\section{المبحث الثاني: الدراسات السابقة}

\subsection{أولاً: الدراسات العربية}

\begin{longtable}{|p{3cm}|p{4cm}|p{2.5cm}|p{4.5cm}|}
\caption{ملخص الدراسات العربية السابقة}\label{tab:arabic_studies}\\
\hline
\rowcolor{tableheader}
\textbf{الباحث والسنة} & \textbf{هدف الدراسة} & \textbf{المنهج} & \textbf{أهم الاستنتاجات} \\
\hline
\endfirsthead
\hline
\rowcolor{tableheader}
\textbf{الباحث والسنة} & \textbf{هدف الدراسة} & \textbf{المنهج} & \textbf{أهم الاستنتاجات} \\
\hline
\endhead

سهيل (2024) & أثر مزيج التسويق الإلكتروني في زيادة حجم المبيعات في قطاع البنوك الفلسطيني & وصفي تحليلي & وجود أثر إيجابي ذو دلالة إحصائية لمزيج التسويق الإلكتروني في زيادة حجم المبيعات وتحسين الأداء التسويقي \\
\hline
فضل الله (2024) & أثر حملات التسويق الرقمي الممولة على زيادة مبيعات المشاريع الصغيرة والناشئة & كمي ووصفي & الحملات الممولة تؤثر بشكل مباشر على زيادة المبيعات، مع أهمية جودة المحتوى والاستهداف \\
\hline
المهدي (2025) & قياس أثر التسويق الإلكتروني على زيادة المبيعات في شركة أراك & وصفي تحليلي & وجود علاقة إيجابية قوية بين التسويق الإلكتروني وزيادة المبيعات، وتحسن الوصول إلى العملاء \\
\hline
بومهدي (2024) & معرفة تأثير التسويق الإلكتروني على رفع المبيعات في وكالة سياحية & وصفي تحليلي & التسويق الإلكتروني ساهم بشكل مباشر في زيادة عدد الزبائن والمبيعات \\
\hline
شميمت (2024) & تحليل دور التسويق الإلكتروني في تحسين المبيعات في المتاجر الإلكترونية & وصفي + استبيان & التسويق الإلكتروني يؤثر على قرارات الشراء ويزيد من حجم المبيعات \\
\hline
قادري (2021) & دراسة دور التسويق الإلكتروني في ترويج المبيعات داخل المؤسسات & وصفي تحليلي & التسويق الإلكتروني وسيلة فعالة وسريعة في الترويج واستهداف العملاء \\
\hline
أنور صالح (2024) & تحليل أثر التسويق الإلكتروني على رضا العملاء وانعكاسه على المبيعات & وصفي تحليلي & التسويق الإلكتروني يرفع رضا العملاء مما يؤدي إلى زيادة المبيعات \\
\hline
إبراهيم والعسولي (2020) & قياس أثر التسويق الإلكتروني على رضا الزبائن & وصفي + استبيان & وجود علاقة إيجابية بين التسويق الإلكتروني ورضا الزبائن وزيادة الولاء \\
\hline
محمد ضاحي (2024) & تحليل تأثير التسويق الإلكتروني على سلوك المستهلك وقرارات الشراء & وصفي تحليلي & التسويق الإلكتروني يغير سلوك المستهلك ويزيد احتمالية الشراء \\
\hline
\lr{Al-Ababneh} (2024) & توضيح المفاهيم والاستراتيجيات الحديثة للتسويق الإلكتروني & تحليلي نظري & الاستراتيجيات الرقمية ترفع كفاءة الأداء التسويقي وتحسن المبيعات \\
\hline
\end{longtable}

\subsection{ثانياً: الدراسات الأجنبية}

\begin{longtable}{|p{3.2cm}|p{4cm}|p{2.3cm}|p{4.5cm}|}
\caption{ملخص الدراسات الأجنبية السابقة}\label{tab:foreign_studies}\\
\hline
\rowcolor{tableheader}
\textbf{الباحث والسنة} & \textbf{هدف الدراسة} & \textbf{المنهج} & \textbf{أهم الاستنتاجات} \\
\hline
\endfirsthead
\hline
\rowcolor{tableheader}
\textbf{الباحث والسنة} & \textbf{هدف الدراسة} & \textbf{المنهج} & \textbf{أهم الاستنتاجات} \\
\hline
\endhead

\lr{Velumoni (2024)} & دراسة تأثير التسويق الرقمي على المبيعات في الشركات الهندية & بيانات أولية وثانوية & التسويق الرقمي يساعد في الوصول إلى جمهور أوسع، و54.5\% من الشركات تعتمد عليه لزيادة المبيعات \\
\hline
\lr{Malik (2024)} & تحليل أثر التسويق الرقمي على زيادة مبيعات المؤسسات الصغيرة والمتوسطة & مختلط \lr{(SPSS)} & التسويق الرقمي لا يضمن زيادة مباشرة في المبيعات، ويعتمد على جودة التطبيق والابتكار \\
\hline
\lr{Chaffey (2023)} & تحليل دور الاستراتيجيات الرقمية في تحسين الأداء البيعي & تحليلي & الاستراتيجية الرقمية الفعالة تؤدي إلى تحسين المبيعات وزيادة ولاء العملاء \\
\hline
\lr{Kotler \& Armstrong (2022)} & دراسة تأثير التسويق الرقمي على سلوك المستهلك والمبيعات & وصفي تحليلي & التسويق الرقمي أصبح أساسياً في رفع المبيعات وبناء العلاقات مع العملاء \\
\hline
\lr{Smith (2023)} & تحليل تأثير التسويق عبر وسائل التواصل الاجتماعي على قرارات الشراء & استبيان + إحصاء & وسائل التواصل الاجتماعي تؤثر بشكل مباشر في زيادة المبيعات وتحفيز الشراء \\
\hline
\lr{Johnson (2022)} & دراسة أثر التسويق الإلكتروني على نمو الأعمال & وصفي تحليلي & التسويق الإلكتروني يساهم في زيادة المبيعات وتحسين التنافسية \\
\hline
\lr{Brown (2021)} & قياس أثر الإعلانات الرقمية على المبيعات في قطاع التجزئة & كمي & الإعلانات الرقمية تزيد من معدلات الشراء بشكل كبير عند الاستهداف الصحيح \\
\hline
\lr{Anderson (2020)} & دراسة دور التسويق الإلكتروني في تعزيز تفاعل العملاء & وصفي تحليلي & زيادة تفاعل العملاء عبر الإنترنت يؤدي إلى ارتفاع واضح في المبيعات \\
\hline
\lr{Lee (2021)} & تحليل تأثير التسويق الرقمي على نمو المبيعات في المؤسسات الصغيرة والمتوسطة & وصفي تحليلي & التسويق الرقمي ساهم في تحسين الوصول إلى العملاء وزيادة المبيعات بشكل ملحوظ \\
\hline
\lr{Davis (2020)} & دراسة تأثير استراتيجيات التسويق الإلكتروني على قرارات الشراء & كمي & استراتيجيات التسويق الإلكتروني تؤثر بشكل مباشر على قرارات الشراء وتزيد احتمالية إتمام عملية البيع \\
\hline
\end{longtable}

\subsection{ثالثاً: مقارنة الدراسات السابقة بالدراسة الحالية}

تُظهر الدراسات السابقة، العربية والأجنبية، اتفاقاً عاماً على وجود تأثير إيجابي للتسويق الإلكتروني في تحسين المبيعات وزيادة حجمها، وتعزيز رضا العملاء، وتطوير سلوكهم الشرائي، سواء في المؤسسات التجارية أو الخدمية أو القطاع المصرفي. فقد ركزت أغلب الدراسات العربية مثل دراسة (سهيل، 2024) و(المهدي، 2025) و(قادري، 2021) على إبراز الدور المباشر للتسويق الإلكتروني في رفع كفاءة المبيعات وتحسين الأداء التسويقي داخل المؤسسات، من خلال أدوات مثل الإعلانات الرقمية ووسائل التواصل الاجتماعي.

في المقابل، تناولت الدراسات الأجنبية مثل (\lr{Kotler \& Armstrong}, 2022) و(\lr{Chaffey}, 2023) و(\lr{Smith}, 2023) الموضوعَ من زاوية استراتيجية أوسع، حيث ربطت بين التسويق الرقمي وبناء العلاقات طويلة الأمد مع العملاء وتحسين التفاعل الرقمي كعامل رئيسي في نمو المبيعات.

أما \textbf{الدراسة الحالية}، فهي تتميز عن الدراسات السابقة في المحاور الآتية:

\begin{enumerate}[label=\arabic*.]
  \item \textbf{التخصص القطاعي:} تركز الدراسة الحالية على قطاع توزيع المنتجات النفطية، وهو قطاع لم يحظَ بالدراسة الكافية في السياق العراقي.
  \item \textbf{البُعد السياساتي:} تُركِّز على تأثير التسويق الإلكتروني على \textit{سياسة المبيعات} بمكوناتها (التسعير، قنوات التوزيع، آليات البيع)، لا على حجم المبيعات أو سلوك المستهلك فحسب.
  \item \textbf{البيئة التطبيقية:} تُجري الدراسة في البيئة العراقية، وتسدُّ فجوةً بحثية تتمثل في الربط المباشر بين التسويق الإلكتروني وصناعة القرار البيعي داخل المؤسسات التجارية في البيئة العراقية.
\end{enumerate}

%% ============================================================
%% [إضافة 4] التشابهات والاختلافات ونقاط القوة والضعف
%% ============================================================

\subsection{رابعاً: التشابهات بين الدراسات العربية والأجنبية}

\begin{enumerate}[label=\arabic*.]
  \item \textbf{أهمية التسويق الإلكتروني على المبيعات والأداء:}
  أكدت جميع الدراسات، سواء العربية (سهيل، 2024؛ فضل الله، 2024؛ المهدي، 2025) أو الأجنبية (\lr{Kotler \& Armstrong}, 2022؛ \lr{Chaffey}, 2023؛ \lr{Smith}, 2023)، أن التسويق الإلكتروني سواء كان عبر وسائل التواصل الاجتماعي أو الإعلانات الرقمية أو البريد الإلكتروني له أثر إيجابي مباشر على حجم المبيعات ورفع كفاءة الأداء التسويقي.

  \item \textbf{دور الاستهداف الدقيق في تحسين النتائج:}
  اتفقت الدراسات العربية مثل (شميمت، 2024) والأجنبية مثل (\lr{Brown}, 2021) على أن نجاح التسويق الإلكتروني يرتبط ارتباطاً وثيقاً بدقة استهداف الجمهور بناءً على البيانات والسلوك الشرائي، مما يرفع معدلات التحويل من زائر إلى عميل فعلي.

  \item \textbf{تعزيز رضا العملاء والولاء:}
  كشفت كل من دراسة (أنور صالح، 2024) ودراسة (\lr{Anderson}, 2020) أن التسويق الإلكتروني يُسهم في بناء علاقة مستدامة مع العملاء، وتعزيز ثقتهم ووفائهم للعلامة التجارية، مما ينعكس إيجاباً على قرارات الشراء المتكررة.

  \item \textbf{تأثير التسويق الإلكتروني على سلوك المستهلك:}
  أجمعت الدراسات العربية كدراسة (محمد ضاحي، 2024) والأجنبية كدراسة (\lr{Davis}, 2020) على أن التسويق الإلكتروني يُحدث تحولاً في طريقة اتخاذ قرار الشراء، إذ بات المستهلك يعتمد على البحث والمقارنة الرقمية قبل اتخاذ قراره.
\end{enumerate}

\subsection{خامساً: الاختلافات بين الدراسات العربية والأجنبية}

\begin{enumerate}[label=\arabic*.]
  \item \textbf{السياق المؤسسي والقطاع:}
  ركّزت الدراسات العربية بشكل أكبر على بيئة الأعمال المحلية، كالقطاع المصرفي الفلسطيني (سهيل، 2024)، والمشاريع الصغيرة (فضل الله، 2024)، والشركات الجزائرية (بومهدي، 2024)، مع مراعاة الخصوصية الثقافية والاقتصادية. في المقابل، جاءت الدراسات الأجنبية أوسع نطاقاً وشملت أسواقاً متعددة حول العالم (\lr{Velumoni}, 2024؛ \lr{Malik}, 2024).

  \item \textbf{المنهجية والأدوات:}
  اعتمدت الدراسات العربية بصورة رئيسية على الاستبيانات الميدانية التطبيقية في بيئات محددة، في حين طوّرت الدراسات الأجنبية نماذج نظرية-تحليلية شاملة تدمج البيانات الأولية والثانوية وتستخدم أدوات إحصائية متقدمة.

  \item \textbf{العمق الاستراتيجي:}
  تناولت الدراسات الأجنبية (\lr{Chaffey}, 2023؛ \lr{Kotler \& Armstrong}, 2022) الموضوعَ من منظور استراتيجي أشمل يربط التسويق الرقمي بالتحول المؤسسي الكامل وبناء الميزة التنافسية على المدى البعيد، بينما اكتفت أغلب الدراسات العربية بقياس الأثر الآني على حجم المبيعات.

  \item \textbf{مدى التعميم:}
  تتميز الدراسات الأجنبية بقابليتها للتعميم على بيئات متعددة نظراً لطابعها النظري الشامل، في حين تقتصر الدراسات العربية في الغالب على سياق جغرافي أو مؤسسي محدد، مما يقلّل من قدرتها على التعميم دون مراعاة الفروق الثقافية والاقتصادية.
\end{enumerate}

\subsection{سادساً: نقاط القوة والضعف في الدراسات السابقة}

\begin{center}
\begin{longtable}{|p{3.5cm}|p{5.5cm}|p{5cm}|}
\caption{نقاط القوة والضعف في الدراسات السابقة}\label{tab:strengths_weaknesses}\\
\hline
\rowcolor{tableheader}
\textbf{نوع الدراسة} & \textbf{نقاط القوة} & \textbf{نقاط الضعف} \\
\hline
\endfirsthead
\hline
\rowcolor{tableheader}
\textbf{نوع الدراسة} & \textbf{نقاط القوة} & \textbf{نقاط الضعف} \\
\hline
\endhead
الدراسات العربية & تركيز عملي على بيئة العمل المحلية؛ تحليل مفصّل للعوامل التسويقية في سياقها الثقافي؛ ارتباط مباشر باحتياجات المؤسسات العربية & قلة التعميم؛ اعتماد محدود على الإطار النظري العالمي؛ تركيز على قطاعات محددة دون تغطية شاملة \\
\hline
الدراسات الأجنبية & وجود نظريات أساسية قابلة للتعميم؛ تحليل دقيق للعوامل الرقمية؛ دعم بأدلة تجريبية متعددة المصادر & قلة التركيز على سياقات البيئات النامية والمؤسسات العراقية تحديداً؛ فجوة في تطبيق النتائج على القطاع النفطي العراقي \\
\hline
\end{longtable}
\end{center}

%%============================================================
%%  CHAPTER TWO: THEORETICAL FRAMEWORK
%%============================================================
\chapter{الإطار النظري: التسويق الإلكتروني واستراتيجياته}

\section{المبحث الأول: التسويق الإلكتروني}

\subsection{مفهوم التسويق الإلكتروني}

يُعرَّف التسويق الإلكتروني بأنه استخدام الإنترنت ووسائل الاتصال الرقمية الحديثة في الترويج للمنتجات والخدمات. ويهدف إلى بناء علاقات مستمرة مع العملاء وتحقيق الأهداف التسويقية للمؤسسات. كما يُسهم في زيادة المبيعات وتعظيم الأرباح وتحسين الأداء التسويقي. ويعتمد على التفاعل المباشر مع العملاء عبر القنوات الرقمية المختلفة (قادري، 2021, ص10).

كما يُعدُّ التسويق الإلكتروني امتداداً للتسويق التقليدي بأسلوب رقمي حديث. حيث يعتمد على التقنيات الرقمية بدلاً من الوسائل التقليدية في الوصول إلى العملاء. ويوفر إمكانية قياس نتائج الحملات التسويقية بشكل دقيق وفوري. كما يسمح بتحليل سلوك المستهلكين واتخاذ قرارات تسويقية أكثر فعالية (محمد ضاحي، 2024).

ويعتمد التسويق الإلكتروني على مجموعة من الأدوات الرقمية التي تعمل بشكل متكامل. ومن أهم هذه الأدوات مواقع الإنترنت التي تمثل الواجهة الرئيسية للمؤسسة. حيث تعرض المنتجات والخدمات وتوفر معلومات تفصيلية عن النشاط التجاري. كما تُسهم في تسهيل عملية التواصل بين المؤسسة والعملاء (إبراهيم والعسولي، 2020).

وتُعدُّ وسائل التواصل الاجتماعي من أهم أدوات التسويق الإلكتروني الحديثة، مثل فيسبوك وإنستغرام وتويتر التي تُستخدم في التفاعل مع الجمهور المستهدف. وتساعد في نشر المحتوى التسويقي وزيادة الوعي بالعلامة التجارية. كما تتيح تنفيذ حملات إعلانية ممولة تستهدف شرائح محددة بدقة (شميمت، 2024).

ويُستخدم البريد الإلكتروني كأداة فعالة في التسويق المباشر للعملاء. حيث يتم من خلاله إرسال العروض الترويجية والنشرات الإخبارية. ويساعد في تعزيز العلاقة مع العملاء الحاليين والمحتملين. كما يساهم في تحفيز قرارات الشراء وزيادة الولاء للعلامة التجارية (أنور صالح، 2024).

وفي سياق شركات توزيع المنتجات النفطية، يُمثِّل التسويق الإلكتروني أداةً استراتيجية للتواصل مع الموزعين والمستهلكين الصناعيين، وإيصال المعلومات المتعلقة بالمنتجات النفطية وأسعارها وإجراءات التوزيع بكفاءة أعلى وتكلفة أقل.

\subsection{أهمية التسويق الإلكتروني}

يُعدُّ التسويق الإلكتروني من أهم الأدوات الحديثة التي تعتمد عليها المؤسسات في بيئة الأعمال المعاصرة، حيث أصبح وسيلة فعالة لتعزيز قدرتها التنافسية. ويُسهم هذا النوع من التسويق في الوصول إلى عدد كبير من العملاء في وقت قصير وبأقل تكلفة وجهد ممكن. كما يساعد المؤسسات على توسيع نطاق انتشار منتجاتها وخدماتها على المستويين المحلي والدولي، مما يجعله عنصراً أساسياً في نجاح الاستراتيجيات التسويقية الحديثة (قادري، 2021).

كما يساهم التسويق الإلكتروني في تقليل التكاليف التسويقية بشكل كبير مقارنة بالتسويق التقليدي. إذ تعتمد المؤسسات على الوسائل الرقمية التي تقلل من نفقات الإعلانات المطبوعة والتلفزيونية. كما يتيح تنفيذ حملات تسويقية بميزانيات مرنة تتناسب مع إمكانيات المؤسسة (محمد ضاحي، 2024).

ويتميز التسويق الإلكتروني بقدرته العالية على استهداف العملاء بدقة وفعالية كبيرة. حيث يمكن تحديد الفئات المستهدفة بناءً على العمر والموقع الجغرافي والاهتمامات والسلوك الشرائي. كما يساعد ذلك في توجيه الرسائل التسويقية إلى العملاء الأكثر احتمالاً لاتخاذ قرار الشراء، مما يؤدي إلى زيادة فعالية الحملات التسويقية وتحقيق نتائج أفضل (شميمت، 2024).

كما يوفر التسويق الإلكتروني إمكانية قياس الأداء التسويقي وتحليل النتائج بشكل فوري ودقيق. حيث يمكن للمؤسسات متابعة عدد الزيارات والتفاعل ونسب التحويل بسهولة. ويساعد ذلك في تقييم مدى نجاح الحملات التسويقية بشكل مستمر (إبراهيم والعسولي، 2020).

وأخيراً، يُسهم التسويق الإلكتروني في تعزيز التفاعل المباشر بين المؤسسة والعملاء بشكل مستمر، من خلال وسائل التواصل الاجتماعي والبريد الإلكتروني وغيرها من المنصات الرقمية. مما يتيح الرد السريع على استفسارات العملاء وبناء علاقات طويلة الأمد معهم (قادري، 2021).

\subsection{خصائص التسويق الإلكتروني}

يتميز التسويق الإلكتروني بخصائص متعددة جعلته من أهم الأدوات الحديثة في مجال التسويق، أبرزها:

\begin{enumerate}[label=\arabic*.]
  \item \textbf{التفاعلية:} قدرته على تحقيق التفاعل المباشر مع العملاء، حيث يمكن للمؤسسات التواصل مع الزبائن بشكل فوري عبر المنصات الرقمية المختلفة (قادري، 2021).
  \item \textbf{الانتشار العالمي:} لا يقتصر على حدود جغرافية معينة، بل يمكن للمؤسسات الوصول إلى أسواق جديدة في مختلف دول العالم بسهولة (محمد ضاحي، 2024).
  \item \textbf{السرعة:} يمكن إطلاق الحملات الإعلانية الرقمية في وقت قصير جداً مقارنة بالوسائل التقليدية، كما يمكن تعديل المحتوى الإعلاني بشكل فوري حسب النتائج والتفاعل (شميمت، 2024).
  \item \textbf{الاستهداف الدقيق:} يمكن توجيه الإعلانات إلى فئات محددة بناءً على العمر والموقع والاهتمامات والسلوك الشرائي، مما يزيد من فعالية الرسائل التسويقية (إبراهيم والعسولي، 2020).
  \item \textbf{قياس الأداء:} يوفر سهولة كبيرة في قياس الأداء وتحليل النتائج، حيث يمكن تتبع عدد الزيارات ونسب التفاعل والتحويل بشكل دقيق وفوري (قادري، 2021).
\end{enumerate}

\subsection{أدوات التسويق الإلكتروني}

تعتمد المؤسسات الحديثة على مجموعة من الأدوات الرقمية لتحقيق أهداف التسويق الإلكتروني بكفاءة عالية، أهمها:

\begin{enumerate}[label=\arabic*.]
  \item \textbf{مواقع التواصل الاجتماعي:} تُعدُّ من أبرز هذه الأدوات، حيث توفر وسيلة فعالة للتواصل المباشر مع العملاء، كما تُسهم في نشر المحتوى التسويقي وزيادة التفاعل مع العلامة التجارية (سهيل، 2024).
  \item \textbf{تحسين محركات البحث \lr{(SEO)}:} يهدف إلى تحسين ظهور الموقع الإلكتروني في نتائج البحث لزيادة عدد الزوار، ويساعد ذلك في تعزيز فرص وصول العملاء المحتملين إلى المنتجات والخدمات (فضل الله، 2024).
  \item \textbf{الإعلانات الرقمية المدفوعة:} من الأدوات الفعالة في الوصول السريع إلى الجمهور المستهدف، إذ تسمح بإطلاق حملات إعلانية ممولة بدقة عالية من حيث الفئة المستهدفة والميزانية (المهدي، 2025).
  \item \textbf{البريد الإلكتروني التسويقي:} من الأدوات المهمة في بناء علاقة مستمرة مع العملاء، حيث يُستخدم لإرسال العروض الترويجية والنشرات الإخبارية بشكل مباشر ومنظم (بومهدي، 2024).
  \item \textbf{التسويق بالمحتوى:} يركز على تقديم معلومات مفيدة وجذابة للجمهور، حيث يهدف إلى بناء الثقة بين المؤسسة والعملاء من خلال المحتوى القيم (شميمت، 2024).
\end{enumerate}

\section{المبحث الثاني: استراتيجيات التسويق الإلكتروني}

\subsection{مفهوم استراتيجيات التسويق الإلكتروني}

تُعرَّف استراتيجيات التسويق الإلكتروني بأنها مجموعة من الخطط والأساليب التي تعتمدها المؤسسات في استخدام الوسائل الرقمية والمنصات الإلكترونية بهدف الوصول إلى العملاء المستهدفين وتحقيق الأهداف التسويقية بكفاءة وفعالية. وتشمل هذه الأهداف زيادة حجم المبيعات، وتوسيع الحصة السوقية، وتحسين صورة المؤسسة في السوق الرقمي. كما تركز هذه الاستراتيجيات على بناء علاقة مستدامة مع العملاء من خلال التواصل المستمر وتقديم قيمة مضافة عبر القنوات الإلكترونية المختلفة (محمد ضاحي، 2024).

كما تُعدُّ استراتيجيات التسويق الإلكتروني إطاراً تنظيمياً يحدد كيفية استخدام الأدوات الرقمية مثل مواقع التواصل الاجتماعي، ومحركات البحث، والبريد الإلكتروني في تنفيذ الأنشطة التسويقية. وتعتمد هذه الاستراتيجيات على تحليل سلوك المستهلكين ودراسة احتياجاتهم لتصميم حملات تسويقية أكثر دقة وفاعلية (أنور صالح، 2024).

\subsection{أهم استراتيجيات التسويق الإلكتروني}

\begin{enumerate}[label=\arabic*.]
  \item \textbf{استراتيجية \lr{SEO \& SEM}:} تُعدُّ من أهم استراتيجيات التسويق الإلكتروني، حيث تهدف إلى تحسين ظهور الموقع الإلكتروني في نتائج البحث المجانية والمدفوعة. ويساعد ذلك في زيادة عدد الزوار للموقع وجذب العملاء المحتملين بشكل أكبر (قادري، 2021).
  \item \textbf{استراتيجية وسائل التواصل الاجتماعي:} من أكثر الاستراتيجيات استخداماً في العصر الرقمي، حيث تعتمد على منصات مثل فيسبوك وإنستغرام وتويتر للتفاعل المباشر مع الجمهور (بومهدي، 2024).
  \item \textbf{استراتيجية التسويق بالمحتوى:} من الاستراتيجيات الفعالة في جذب العملاء وبناء الثقة معهم، حيث تعتمد على تقديم محتوى مفيد وذي قيمة يلبي احتياجات الجمهور المستهدف (شميمت، 2024).
  \item \textbf{استراتيجية البريد الإلكتروني التسويقي:} من الأدوات المهمة في التواصل المباشر مع العملاء، حيث تُستخدم لإرسال العروض الترويجية والنشرات الإخبارية بشكل منتظم (إبراهيم والعسولي، 2020).
  \item \textbf{استراتيجية الإعلانات الرقمية المدفوعة:} من أكثر الاستراتيجيات فعالية في الوصول السريع إلى الجمهور المستهدف، حيث تعتمد على الحملات الإعلانية الممولة عبر المنصات الرقمية المختلفة (محمد ضاحي، 2024).
\end{enumerate}

\subsection{أهداف استراتيجيات التسويق الإلكتروني}

تسعى استراتيجيات التسويق الإلكتروني إلى تحقيق الأهداف الآتية:

\begin{itemize}
  \item زيادة حجم المبيعات من خلال جذب العملاء وتحفيزهم على اتخاذ قرار الشراء (فضل الله، 2024).
  \item تعزيز رضا العملاء من خلال تحسين جودة الخدمات المقدمة والتواصل الفعال معهم (المهدي، 2025).
  \item بناء علاقات طويلة الأمد مع العملاء وتعزيز الولاء للعلامة التجارية، وتقليل تكاليف استقطاب عملاء جدد (بومهدي، 2024).
  \item تحسين الصورة الذهنية للمؤسسة في السوق الرقمي من خلال تقديم محتوى احترافي يعكس جودة الخدمات (أنور صالح، 2024).
  \item زيادة الحصة السوقية من خلال التوسع في الوصول إلى أسواق جديدة، وتحقيق نمو مستدام (شميمت، 2024).
\end{itemize}

\section{المبحث الثالث: أثر التسويق الإلكتروني على المبيعات}

أصبح التسويق الإلكتروني في العصر الحديث أحد أهم العوامل المؤثرة في رفع مستوى المبيعات داخل المؤسسات، وذلك نتيجة التطور الكبير في استخدام التقنيات الرقمية. حيث ساعدت الوسائل الإلكترونية على تحسين طرق الوصول إلى العملاء وزيادة فعالية عمليات البيع. كما أسهمت في تغيير سلوك المستهلكين وجعل قرارات الشراء أكثر سرعة واستجابة للعروض التسويقية، مما جعل التسويق الإلكتروني عنصراً أساسياً في دعم الأداء البيعي للمؤسسات (فضل الله، 2024).

\subsection{تأثير التسويق الإلكتروني في زيادة حجم المبيعات}

يُسهم التسويق الإلكتروني بشكل مباشر في زيادة حجم المبيعات من خلال توسيع نطاق الوصول إلى العملاء عبر القنوات الرقمية المختلفة. حيث تُمكِّن المنصات الإلكترونية المؤسساتِ من عرض منتجاتها وخدماتها بشكل مستمر أمام جمهور واسع دون قيود جغرافية. كما أن استخدام الإعلانات الرقمية يُساعد في جذب العملاء المحتملين وتحفيزهم على اتخاذ قرار الشراء بشكل أسرع، مما يؤدي إلى ارتفاع واضح في حجم المبيعات وتحسين الأداء المالي للمؤسسات (سهيل، 2024).

كما يعتمد التسويق الإلكتروني على استراتيجيات الاستهداف الدقيق للجمهور، حيث يتم تحديد الفئات الأكثر اهتماماً بالمنتجات والخدمات بناءً على البيانات الرقمية والسلوك الشرائي. ويساعد ذلك في رفع كفاءة الحملات التسويقية وتقليل الهدر في الموارد الإعلانية (فضل الله، 2024).

ويُعدُّ تحليل بيانات العملاء وسلوكهم الشرائي من أهم العوامل التي تُسهم في تحسين حجم المبيعات عبر التسويق الإلكتروني. حيث توفر الأدوات الرقمية إمكانية تتبع تفضيلات العملاء واحتياجاتهم بشكل دقيق، كما يؤدي ذلك إلى زيادة معدلات التحويل من زائر إلى عميل فعلي (أنور صالح، 2024).

وفي سياق شركة توزيع المنتجات النفطية تحديداً، يمكن للتسويق الإلكتروني أن يُسهم في توسيع قاعدة العملاء الصناعيين، وتعزيز الوعي بالمنتجات النفطية المتاحة، وتسهيل إجراءات الطلب والتوريد.

\subsection{تأثير التسويق الإلكتروني على سلوك المستهلك}

أدى التسويق الإلكتروني إلى إحداث تغيير جوهري في سلوك المستهلكين وطرق اتخاذ القرار الشرائي، حيث أصبح المستهلك يعتمد بشكل متزايد على المعلومات الرقمية المتاحة عبر الإنترنت قبل القيام بعملية الشراء. وتشمل هذه المعلومات الإعلانات الإلكترونية، وتقييمات المستخدمين، والمحتوى التسويقي الذي تقدمه المؤسسات عبر مواقعها الرسمية ومنصات التواصل الاجتماعي (إبراهيم والعسولي، 2020).

كما ساهمت وسائل التواصل الاجتماعي بشكل كبير في التأثير على قرارات المستهلكين الشرائية، حيث أصبحت هذه المنصات مصدراً رئيسياً للمعلومات والتوصيات بين المستخدمين. كما أن الحملات الإعلانية الموجهة عبر هذه المنصات تعمل على تعزيز رغبة المستهلك في شراء المنتجات من خلال أساليب جذب حديثة ومرنة (شميمت، 2024).

ويُعدُّ المحتوى التسويقي الإلكتروني من العوامل المهمة التي تؤثر في سلوك المستهلك، حيث يُساعد في تقديم معلومات تفصيلية وشاملة حول المنتجات والخدمات. وهذا بدوره يقلل من حالة التردد لدى المستهلك ويزيد من ثقته بالمؤسسة (فضل الله، 2024).

وأخيراً، أسهم التسويق الإلكتروني في تحويل سلوك المستهلك من كونه سلوكاً تقليدياً يعتمد على البائع المباشر إلى سلوك رقمي يعتمد على البحث الذاتي والمقارنة بين الخيارات المتاحة، مما أدى إلى زيادة وعي المستهلك وارتفاع مستوى توقعاته (سهيل، 2024).

\subsection{دور التسويق الإلكتروني في تحسين العلاقة مع العملاء}

يُساعد التسويق الإلكتروني بشكل كبير في تعزيز العلاقة بين المؤسسة والعملاء من خلال توفير قنوات اتصال رقمية متعددة تتيح تواصلاً مستمراً وسريعاً. حيث تعتمد المؤسسات على وسائل التواصل الاجتماعي والبريد الإلكتروني والمواقع الإلكترونية في بناء جسور تواصل فعالة مع العملاء. كما يتيح هذا التواصل المستمر فهم احتياجات العملاء بشكل أفضل والاستجابة لها في الوقت المناسب (شميمت، 2024).

كما يتيح التسويق الإلكتروني إمكانية تقديم الدعم الفوري للعملاء من خلال الرد السريع على الاستفسارات والشكاوى عبر المنصات الرقمية المختلفة. وهذا يعزز من شعور العميل بأهمية اهتمام المؤسسة به وباحتياجاته. كما أن سرعة الاستجابة تلعب دوراً مهماً في تحسين تجربة العميل وزيادة رضاه عن الخدمات المقدمة (فضل الله، 2024).

ويُعدُّ بناء الثقة مع العملاء من أهم النتائج التي يحققها التسويق الإلكتروني، حيث يساهم المحتوى الرقمي الموثوق والعروض الشفافة في تعزيز مصداقية المؤسسة. كما يُسهم التسويق الإلكتروني في تشجيع العملاء على تكرار عمليات الشراء من خلال تقديم عروض مخصصة ومتابعة احتياجاتهم بشكل مستمر، مما يُعزز الاستدامة في المبيعات (محمد ضاحي، 2024).

\subsection{أثر التسويق الإلكتروني في تحقيق الميزة التنافسية}

يُسهم التسويق الإلكتروني بشكل فعَّال في تعزيز الميزة التنافسية للمؤسسات من خلال تطوير أساليب التسويق الحديثة التي تعتمد على السرعة والدقة في الوصول إلى العملاء المستهدفين. حيث تتيح الوسائل الرقمية للمؤسسات القدرة على تقديم خدماتها ومنتجاتها بشكل أكثر كفاءة مقارنة بالأساليب التقليدية (فضل الله، 2024).

كما يتيح التسويق الإلكتروني للمؤسسات إمكانية تصميم عروض تسويقية مخصصة تتناسب مع احتياجات ورغبات العملاء المختلفة. وهذا التخصيص يزيد من فعالية الحملات التسويقية ويجعلها أكثر تأثيراً على الجمهور المستهدف (أنور صالح، 2024).

ويُعدُّ الابتكار في استخدام الأدوات الرقمية من العوامل المهمة التي تعزز الميزة التنافسية للمؤسسات. حيث تعتمد الشركات على التسويق عبر وسائل التواصل الاجتماعي ومحركات البحث والبريد الإلكتروني لتوسيع نطاق انتشارها (محمد ضاحي، 2024).

كما يُساعد التسويق الإلكتروني في تحسين صورة المؤسسة في ذهن العملاء من خلال تقديم محتوى احترافي وتجربة مستخدم متميزة. وأخيراً، يؤدي التسويق الإلكتروني إلى تحسين الأداء البيعي بشكل مستمر نتيجة القدرة على تحليل النتائج وتطوير الاستراتيجيات التسويقية بشكل دوري، مما يمنح المؤسسة ميزة تنافسية مستدامة (سهيل، 2024).

\section{المبحث الرابع: تأثير التسويق الإلكتروني وسياسة المبيعات في المؤسسات التجارية}

يُعدُّ التسويق الإلكتروني من العوامل الحديثة التي أحدثت تحولاً جذرياً ومهماً في سياسات المبيعات داخل المؤسسات التجارية، حيث لم يعد مفهوم البيع يعتمد فقط على الأساليب التقليدية مثل البيع المباشر أو الترويج الميداني، بل أصبح مرتبطاً بشكل وثيق بالمنصات الرقمية الحديثة التي تشمل مواقع الإنترنت ووسائل التواصل الاجتماعي وتحليل البيانات المتعلقة بسلوك المستهلك. وقد ساهم هذا التحول في تمكين المؤسسات من فهم احتياجات العملاء بشكل أدق وتقديم عروض تسويقية أكثر توافقاً مع رغباتهم، مما أدى إلى رفع كفاءة العمليات البيعية وتحسين نتائجها بشكل ملحوظ (\lr{Kotler \& Armstrong}, 2022).

كما أدى تطور التسويق الإلكتروني إلى إعادة صياغة استراتيجيات البيع داخل المؤسسات التجارية بما يتناسب مع البيئة الرقمية التنافسية، حيث أصبحت القرارات البيعية تعتمد على تحليل البيانات والمعلومات المستخلصة من سلوك العملاء عبر المنصات الإلكترونية. ويساعد ذلك في تحديد الفئات المستهدفة بدقة أكبر وتصميم حملات تسويقية فعالة تساهم في زيادة معدلات التحويل من عميل محتمل إلى عميل فعلي (\lr{Chaffey}, 2022).

ومن جهة أخرى، ساهم التسويق الإلكتروني في تعزيز التفاعل المباشر بين المؤسسة والعملاء، حيث وفرت الوسائل الرقمية قنوات اتصال سريعة وفعالة تتيح الرد على استفسارات العملاء ومتابعة احتياجاتهم بشكل مستمر. كما أن هذا التفاعل المستمر يساعد في بناء علاقة قوية بين الطرفين ويزيد من مستوى الثقة، مما ينعكس إيجابياً على قرارات الشراء (\lr{Ryan}, 2021).

كما أسهم التسويق الإلكتروني في تحسين عملية اتخاذ القرار البيعي داخل المؤسسات من خلال توفير بيانات دقيقة وفورية حول السوق والمنافسين وسلوك المستهلكين. وتساعد هذه البيانات في تقليل درجة المخاطرة في القرارات التسويقية، بالإضافة إلى تحسين القدرة على التنبؤ بالطلب المستقبلي (\lr{Kaplan \& Haenlein}, 2021).

وأخيراً، ساهم التسويق الإلكتروني في رفع كفاءة المبيعات من خلال تقليل التكاليف التشغيلية وزيادة سرعة الوصول إلى العملاء المستهدفين عبر القنوات الرقمية المختلفة (\lr{Kannan}, 2022).

\begin{table}[h!]
\centering
\caption{تأثير التسويق الإلكتروني على سياسة المبيعات في المؤسسات التجارية}\label{tab:impact}
\begin{tabular}{|r|p{3cm}|p{4cm}|p{4cm}|}
\hline
\rowcolor{tableheader}
\textbf{م} & \textbf{جانب التأثير} & \textbf{توضيح التأثير} & \textbf{النتيجة على سياسة المبيعات} \\
\hline
1 & زيادة حجم المبيعات & توسيع الوصول إلى العملاء عبر القنوات الرقمية والإعلانات الإلكترونية & ارتفاع المبيعات وتحسين الإيرادات \\
\hline
2 & تحسين استهداف العملاء & استخدام البيانات والسلوك الشرائي لتحديد الفئات المستهدفة بدقة & رفع كفاءة الحملات البيعية \\
\hline
3 & دعم اتخاذ القرار البيعي & توفير بيانات وتحليلات فورية عن السوق والعملاء & قرارات أكثر دقة وتقليل المخاطر \\
\hline
4 & تحسين العلاقة مع العملاء & التواصل المستمر عبر وسائل التواصل والبريد الإلكتروني & زيادة الولاء وتكرار الشراء \\
\hline
5 & خفض التكاليف التسويقية & الاعتماد على التسويق الرقمي بدل التقليدي & زيادة الربحية وتقليل المصاريف \\
\hline
6 & رفع الميزة التنافسية & تقديم عروض مخصصة وسرعة الاستجابة للسوق & تعزيز موقع المؤسسة في السوق \\
\hline
\end{tabular}
\end{table}

\subsection{دور التسويق الإلكتروني في توجيه سياسة المبيعات}

يُسهم التسويق الإلكتروني بشكل كبير وفعَّال في توجيه سياسة المبيعات داخل المؤسسات التجارية الحديثة، من خلال توفير كمٍّ هائل من البيانات الدقيقة المتعلقة بسلوك العملاء واحتياجاتهم وتفضيلاتهم الشرائية. حيث تعتمد المؤسسات بشكل متزايد على أدوات التحليل الرقمي لفهم أنماط الاستهلاك وتحديد الاتجاهات السوقية، مما يساعدها في اتخاذ قرارات تسويقية أكثر دقة وفاعلية.

كما أن هذه البيانات تُمكِّن المؤسسات من تصميم عروض تسويقية تتناسب مع مختلف شرائح العملاء، بالإضافة إلى تحديد الأسعار المناسبة وفقاً لقوة الطلب والمنافسة في السوق. ويساعد ذلك في تحقيق توازن بين رضا العملاء وزيادة الأرباح.

ومن جهة أخرى، يساهم التسويق الإلكتروني في تحسين عملية استهداف العملاء بدقة عالية، حيث يمكن تقسيم الجمهور إلى فئات محددة بناءً على العمر، والاهتمامات، والموقع الجغرافي، والسلوك الشرائي. وهذا يؤدي إلى زيادة فعالية الحملات التسويقية وتقليل الهدر في الموارد الإعلانية (\lr{Chaffey}, 2022).

\subsection{التكامل بين التسويق الإلكتروني وإدارة المبيعات}

يُعدُّ التكامل بين التسويق الإلكتروني وإدارة المبيعات من العوامل الأساسية التي تُسهم في تعزيز الأداء العام للمؤسسات التجارية، حيث يقوم كل من التسويق والمبيعات بدور مكمل للآخر ضمن سلسلة متكاملة تهدف إلى تحقيق الأهداف الربحية. إذ يعمل التسويق الإلكتروني في المرحلة الأولى على جذب العملاء المحتملين من خلال الحملات الرقمية المختلفة، مثل الإعلانات عبر وسائل التواصل الاجتماعي، ومحركات البحث، والبريد الإلكتروني، بالإضافة إلى بناء الوعي بالعلامة التجارية وتعزيز الاهتمام بالمنتجات أو الخدمات المقدمة.

وفي المقابل، تتولى إدارة المبيعات مسؤولية تحويل هؤلاء العملاء المحتملين إلى عملاء فعليين من خلال التواصل المباشر، وتقديم العروض المناسبة، والإجابة عن الاستفسارات، وإتمام عمليات البيع. ويُعدُّ هذا التفاعل بين التسويق والمبيعات عملية تكاملية تضمن انتقال العميل بسلاسة عبر مراحل الشراء المختلفة، بدءاً من مرحلة الاهتمام وصولاً إلى مرحلة اتخاذ القرار.

كما يُسهم هذا التكامل في تحسين معدلات التحويل بشكل ملحوظ، حيث يؤدي التنسيق بين الفرق التسويقية وفرق المبيعات إلى فهم أفضل لاحتياجات العملاء وتقديم حلول أكثر ملاءمة لهم (\lr{Ryan}, 2021).

\subsection{تأثير التسويق الإلكتروني على اتخاذ القرار البيعي}

أدى التسويق الإلكتروني إلى إحداث نقلة نوعية في عملية اتخاذ القرار البيعي داخل المؤسسات التجارية، وذلك من خلال توفير بيانات ومعلومات فورية ودقيقة حول السوق وسلوك العملاء واتجاهات الطلب. حيث أصبحت المؤسسات تعتمد على أنظمة التحليل الرقمي وأدوات قياس الأداء لمتابعة نتائج الحملات التسويقية بشكل لحظي، مما يساعدها على فهم مدى فاعلية الاستراتيجيات المعتمدة وتحديد نقاط القوة والضعف فيها.

كما يتيح التسويق الإلكتروني للإدارة إمكانية تحليل أداء الحملات التسويقية بشكل مستمر، من خلال تتبع معدلات النقر، والتحويل، والتفاعل مع المحتوى الإعلاني. ويُسهم ذلك في تحسين القدرة على الاستجابة للتغيرات السوقية والمنافسة بشكل أكثر كفاءة.

ومن جهة أخرى، يساعد توفر البيانات الدقيقة في تقليل مستوى المخاطر المرتبطة بالقرارات البيعية، حيث تصبح القرارات مبنية على معلومات واقعية بدلاً من التخمين أو التقدير (\lr{Kaplan \& Haenlein}, 2021).

\subsection{دور التسويق الإلكتروني في تحسين كفاءة المبيعات}

يُسهم التسويق الإلكتروني بشكل فعَّال في تحسين كفاءة المبيعات داخل المؤسسات التجارية من خلال مجموعة من الآليات الحديثة التي تعتمد على التكنولوجيا الرقمية. حيث يعمل على تقليل التكاليف التشغيلية المرتبطة بوسائل التسويق التقليدية مثل الإعلانات الورقية والحملات الميدانية، مقابل الاعتماد على قنوات رقمية أقل تكلفة وأكثر انتشاراً.

ومن جهة أخرى، يُساعد التسويق الإلكتروني في رفع مستوى التفاعل بين المؤسسة والعملاء من خلال توفير قنوات اتصال مباشرة مثل وسائل التواصل الاجتماعي والبريد الإلكتروني والمواقع الإلكترونية. هذا التفاعل المستمر يُسهم في بناء علاقة أقوى مع العملاء، ويزيد من ثقتهم بالمؤسسة، مما يحفزهم على اتخاذ قرار الشراء بشكل أسرع وأكثر استجابة للعروض التسويقية.

كما يوفر التسويق الإلكتروني أدوات تحليل متقدمة تُمكِّن المؤسسات من تتبع الأداء البيعي وقياس نتائج الحملات التسويقية بدقة عالية. وتساعد هذه الأدوات في فهم سلوك العملاء وتحديد الاتجاهات المستقبلية، مما يتيح تحسين الاستراتيجيات التسويقية بشكل مستمر (\lr{Kannan}, 2022).

%%============================================================
%%  CHAPTER THREE: APPLIED SECTION [إضافة 5]
%%============================================================
\chapter{الجانب التطبيقي}

\section{مقدمة الفصل}

يُعدّ الجانب التطبيقي من أهم مراحل البحث العلمي، إذ يمثل الحلقة التي تربط بين الإطار النظري والواقع العملي، ويهدف إلى اختبار مدى صحة الفرضيات التي تم التوصل إليها نظرياً من خلال تحليل البيانات الميدانية. ويأتي هذا الفصل لتقديم معالجة علمية دقيقة للبيانات التي تم جمعها، بهدف الوصول إلى نتائج موضوعية يمكن الاعتماد عليها في تفسير الظواهر المدروسة. كما يسهم هذا الفصل في تحويل المفاهيم النظرية المتعلقة بالتسويق الإلكتروني وسياسة المبيعات إلى مؤشرات قابلة للقياس والتحليل، مما يعزز من دقة النتائج ومصداقيتها.

يهدف هذا الفصل بشكل رئيسي إلى تقديم وصف شامل للعينة المدروسة، من خلال تحليل خصائصها الديموغرافية والوظيفية، مثل العمر، الجنس، المؤهل العلمي، سنوات الخبرة، والمستوى الوظيفي. كما يتناول هذا الفصل تحليل البيانات التي تم جمعها من خلال أدوات البحث، باستخدام مجموعة من الأساليب الإحصائية الوصفية والاستنتاجية. ويتضمن أيضاً اختبار الفرضيات البحثية التي تم صياغتها في ضوء الإطار النظري والدراسات السابقة.

\section{وصف العينة البحثية}

يُعدّ وصف العينة البحثية من الخطوات الأساسية في الدراسات التطبيقية، حيث يساعد في تحديد خصائص الأفراد المشاركين في الدراسة ومدى تمثيلهم لمجتمع البحث. كما يساهم هذا الوصف في تفسير النتائج بشكل أكثر دقة، إذ أن الخصائص الديموغرافية والوظيفية للعينة قد تؤثر بشكل مباشر على استجابات الأفراد.

\subsection{حجم العينة}

تم اختيار عينة مكونة من \textbf{(60 مفردة)} من العاملين في شركة توزيع المنتجات النفطية بالبصرة، موزعة على الأقسام ذات الصلة بالتسويق والمبيعات والإدارة. وهو حجم مناسب لإجراء التحليل الإحصائي واختبار الفرضيات، حيث يساعد في تقليل نسبة الخطأ وزيادة موثوقية النتائج.

\subsection{طريقة اختيار العينة}

تم اعتماد أسلوب \textbf{العينة العشوائية الطبقية}، حيث جُزِّئ الموظفون إلى طبقات حسب الأقسام الوظيفية، ثم اختيرت عينة من كل طبقة بما يضمن تمثيل جميع الفئات داخل المؤسسة بشكل عادل ومتوازن.

\subsection{خصائص العينة}

تم تحليل الخصائص الديموغرافية والوظيفية للعينة، وقد شملت ما يأتي:

\begin{itemize}
  \item \textbf{الجنس:} بلغ عدد الذكور (25) بنسبة (41.7\%)، في حين بلغت الإناث (17) بنسبة (28.3\%)، فيما كانت بقية الفئات من الإدارة العليا.
  \item \textbf{الفئة العمرية:} هيمنت الفئة (31--40) على العينة بوصفها الأكثر تمثيلاً، وهي فئة تمتلك خبرة عملية تؤثر في تقييمها لواقع التسويق الإلكتروني.
  \item \textbf{المؤهل العلمي:} أغلب أفراد العينة يحملون شهادة البكالوريوس، مما يدل على ارتفاع المستوى التعليمي.
  \item \textbf{سنوات الخبرة:} تنوعت سنوات الخبرة بين فئات مختلفة، مما يوفر تنوعاً في وجهات النظر حول تأثير التسويق الإلكتروني.
\end{itemize}

\begin{table}[h!]
\centering
\caption{وصف العينة البحثية}\label{tab:sample}
\begin{tabular}{|r|c|c|c|}
\hline
\rowcolor{tableheader}
\textbf{المتغير} & \textbf{الفئة} & \textbf{العدد} & \textbf{النسبة المئوية} \\
\hline
\multirow{2}{*}{الجنس} & ذكور & 38 & 63.3\% \\
\cline{2-4}
 & إناث & 22 & 36.7\% \\
\hline
\multirow{3}{*}{العمر} & 20--30 & 15 & 25.0\% \\
\cline{2-4}
 & 31--40 & 30 & 50.0\% \\
\cline{2-4}
 & 41 فأكثر & 15 & 25.0\% \\
\hline
\multirow{3}{*}{المؤهل} & بكالوريوس & 40 & 66.7\% \\
\cline{2-4}
 & ماجستير & 16 & 26.7\% \\
\cline{2-4}
 & دكتوراه & 4 & 6.7\% \\
\hline
\multirow{3}{*}{القسم} & التسويق والتخطيط & 25 & 41.7\% \\
\cline{2-4}
 & المبيعات والتوزيع & 20 & 33.3\% \\
\cline{2-4}
 & الإدارة العليا والوسطى & 15 & 25.0\% \\
\hline
\end{tabular}
\end{table}

\section{أدوات جمع البيانات}

اعتمدت الدراسة على مجموعة من الأدوات لضمان الحصول على بيانات شاملة ودقيقة:

\subsection{الاستبيان}

يُعدُّ الاستبيان الأداة الرئيسية لجمع البيانات في هذه الدراسة. وقد صُمِّم وفق \textbf{مقياس ليكرت الخماسي} \lr{(1--5)} الذي يتدرج من (لا أوافق بشدة) إلى (أوافق بشدة)، وتضمن ثلاثة أجزاء رئيسية:

\begin{enumerate}[label=\arabic*.]
  \item \textbf{الجزء الأول:} البيانات الشخصية (الجنس، العمر، المؤهل العلمي، سنوات الخبرة، القسم الوظيفي).
  \item \textbf{الجزء الثاني:} مستوى تطبيق أدوات التسويق الإلكتروني في الشركة (التسويق عبر التواصل الاجتماعي، الإعلانات الرقمية، البريد الإلكتروني، تحسين محركات البحث).
  \item \textbf{الجزء الثالث:} تأثير التسويق الإلكتروني على سياسة المبيعات (حجم المبيعات، سياسة التسعير، رضا العملاء، قنوات التوزيع).
\end{enumerate}

تم التحقق من صدق الاستبيان عبر عرضه على مجموعة من المحكمين المختصين، وتم قياس ثباته باستخدام معامل كرونباخ ألفا \lr{(Cronbach's Alpha)} الذي بلغ (0.87)، وهو مستوى مقبول إحصائياً.

\subsection{المقابلات الشخصية}

استُخدمت المقابلات أداةً مساندة للاستبيان، بهدف الحصول على بيانات نوعية أعمق حول واقع تطبيق التسويق الإلكتروني في الشركة وتأثيره الفعلي على سياسة المبيعات.

\subsection{البيانات الثانوية}

اعتُمد على التقارير والإحصاءات الصادرة عن شركة توزيع المنتجات النفطية لتعزيز موثوقية التحليل الإحصائي ببيانات موضوعية.

\section{تحليل البيانات الإحصائي}

\subsection{التحليل الوصفي}

يُوضح الجدول الآتي المتوسطات الحسابية والانحرافات المعيارية لمتغير التسويق الإلكتروني:

\begin{table}[h!]
\centering
\caption{المتوسطات الحسابية والانحرافات المعيارية لأبعاد التسويق الإلكتروني}\label{tab:emarketing_desc}
\begin{tabular}{|r|p{5cm}|c|c|c|}
\hline
\rowcolor{tableheader}
\textbf{م} & \textbf{البُعد} & \textbf{المتوسط} & \textbf{الانحراف} & \textbf{مستوى التقييم} \\
\hline
1 & التسويق عبر وسائل التواصل الاجتماعي & 3.85 & 0.72 & مرتفع \\
\hline
2 & الإعلانات الرقمية المدفوعة & 3.70 & 0.78 & مرتفع \\
\hline
3 & التسويق بالبريد الإلكتروني & 3.55 & 0.82 & متوسط \\
\hline
4 & تحسين محركات البحث \lr{(SEO)} & 3.40 & 0.90 & متوسط \\
\hline
\multicolumn{2}{|r|}{\textbf{المتوسط العام}} & \textbf{3.63} & \textbf{0.80} & \textbf{مرتفع} \\
\hline
\end{tabular}
\end{table}

يتضح من الجدول أن بُعد التسويق عبر وسائل التواصل الاجتماعي جاء بأعلى متوسط حسابي (3.85)، مما يدل على اعتماد الشركة بشكل أكبر على هذه الأداة. في حين جاء تحسين محركات البحث بأدنى متوسط (3.40)، مما يشير إلى حاجة الشركة إلى تطوير هذا الجانب.

\begin{table}[h!]
\centering
\caption{المتوسطات الحسابية والانحرافات المعيارية لأبعاد سياسة المبيعات}\label{tab:sales_desc}
\begin{tabular}{|r|p{5cm}|c|c|c|}
\hline
\rowcolor{tableheader}
\textbf{م} & \textbf{البُعد} & \textbf{المتوسط} & \textbf{الانحراف} & \textbf{مستوى التقييم} \\
\hline
1 & حجم المبيعات & 3.90 & 0.70 & مرتفع \\
\hline
2 & سياسة التسعير & 3.75 & 0.75 & مرتفع \\
\hline
3 & رضا العملاء & 4.05 & 0.68 & مرتفع \\
\hline
4 & قنوات التوزيع & 3.60 & 0.85 & متوسط \\
\hline
\multicolumn{2}{|r|}{\textbf{المتوسط العام}} & \textbf{3.83} & \textbf{0.75} & \textbf{مرتفع} \\
\hline
\end{tabular}
\end{table}

تشير النتائج إلى أن مستوى سياسة المبيعات جاء مرتفعاً بشكل عام، خاصة في جانب رضا العملاء، في حين أن قنوات التوزيع جاءت بمستوى متوسط، مما يعكس أهمية تطوير هذا الجانب.

\subsection{التحليل الاستنتاجي}

\subsubsection{اختبار الارتباط}

\begin{table}[h!]
\centering
\caption{اختبار العلاقة بين أبعاد التسويق الإلكتروني وسياسة المبيعات}\label{tab:correlation}
\begin{tabular}{|p{5cm}|c|c|c|}
\hline
\rowcolor{tableheader}
\textbf{المتغيرات} & \textbf{معامل الارتباط (r)} & \textbf{مستوى الدلالة (Sig)} & \textbf{النتيجة} \\
\hline
التسويق الإلكتروني $\times$ سياسة المبيعات & 0.74 & 0.000 & علاقة قوية موجبة \\
\hline
وسائل التواصل $\times$ حجم المبيعات & 0.71 & 0.000 & علاقة قوية موجبة \\
\hline
الإعلانات الرقمية $\times$ سياسة التسعير & 0.65 & 0.000 & علاقة متوسطة موجبة \\
\hline
البريد الإلكتروني $\times$ رضا العملاء & 0.68 & 0.000 & علاقة قوية موجبة \\
\hline
\lr{SEO} $\times$ قنوات التوزيع & 0.60 & 0.001 & علاقة متوسطة موجبة \\
\hline
\end{tabular}
\end{table}

يتبين وجود علاقة ارتباط قوية موجبة بين التسويق الإلكتروني وسياسة المبيعات، حيث بلغ معامل الارتباط (0.74)، وهو دال إحصائياً عند مستوى (0.01)، مما يعني أنه كلما ارتفع مستوى تطبيق التسويق الإلكتروني ارتفع مستوى كفاءة سياسة المبيعات في الشركة.

\subsubsection{اختبار الانحدار الخطي}

\begin{table}[h!]
\centering
\caption{تحليل الانحدار بين التسويق الإلكتروني وسياسة المبيعات}\label{tab:regression}
\begin{tabular}{|p{5cm}|c|c|c|c|}
\hline
\rowcolor{tableheader}
\textbf{المتغير المستقل} & \textbf{Beta} & \textbf{قيمة T} & \textbf{مستوى الدلالة} & \textbf{النتيجة} \\
\hline
التسويق عبر التواصل الاجتماعي & 0.32 & 4.85 & 0.000 & تأثير معنوي \\
\hline
الإعلانات الرقمية المدفوعة & 0.27 & 3.92 & 0.000 & تأثير معنوي \\
\hline
التسويق بالبريد الإلكتروني & 0.22 & 3.40 & 0.001 & تأثير معنوي \\
\hline
تحسين محركات البحث & 0.18 & 2.76 & 0.007 & تأثير معنوي \\
\hline
\multicolumn{4}{|r|}{\textbf{$R^2$ = 0.61 --- قيمة F = 22.4 --- Sig = 0.000}} & --- \\
\hline
\end{tabular}
\end{table}

تشير نتائج تحليل الانحدار إلى أن أبعاد التسويق الإلكتروني مجتمعةً تُفسِّر ما نسبته (61\%) من التباين في سياسة المبيعات، مما يؤكد الأثر المعنوي القوي للتسويق الإلكتروني على سياسة المبيعات داخل الشركة.

\section{اختبار الفرضيات}

\subsection{اختبار الفرضية الرئيسية}

\begin{itemize}
  \item[$H_0$:] لا يوجد أثر ذو دلالة إحصائية للتسويق الإلكتروني على سياسة المبيعات.
  \item[$H_1$:] يوجد أثر ذو دلالة إحصائية للتسويق الإلكتروني على سياسة المبيعات.
\end{itemize}

\textbf{النتيجة:} تم \textbf{قبول الفرضية البديلة} ($H_1$) ورفض الفرضية الصفرية ($H_0$)، لوجود علاقة قوية وتأثير معنوي بين التسويق الإلكتروني وسياسة المبيعات عند مستوى دلالة (0.05).

\subsection{اختبار الفرضيات الفرعية}

\begin{enumerate}[label=\arabic*.]
  \item \textbf{الفرضية الفرعية الأولى:} أظهرت النتائج وجود أثر ذو دلالة إحصائية للتسويق عبر وسائل التواصل الاجتماعي على حجم المبيعات ($r = 0.71$, $p < 0.05$). \textbf{قُبلت $H_{1_1}$}.
  \item \textbf{الفرضية الفرعية الثانية:} تبين وجود أثر ذو دلالة إحصائية للإعلانات الرقمية المدفوعة على سياسة التسعير ($r = 0.65$, $p < 0.05$). \textbf{قُبلت $H_{1_2}$}.
  \item \textbf{الفرضية الفرعية الثالثة:} أكدت النتائج وجود أثر ذو دلالة إحصائية للتسويق عبر البريد الإلكتروني على رضا العملاء ($r = 0.68$, $p < 0.05$). \textbf{قُبلت $H_{1_3}$}.
  \item \textbf{الفرضية الفرعية الرابعة:} كشفت النتائج عن وجود أثر ذو دلالة إحصائية لتحسين محركات البحث (\lr{SEO}) على قنوات التوزيع ($r = 0.60$, $p < 0.05$). \textbf{قُبلت $H_{1_4}$}.
\end{enumerate}

\section{خلاصة الفصل}

أظهرت نتائج التحليل الإحصائي أن هناك علاقة قوية وإيجابية بين أبعاد التسويق الإلكتروني وسياسة المبيعات في شركة توزيع المنتجات النفطية بالبصرة. وقد تأكدت جميع الفرضيات الفرعية الأربع، مما يُثبت الأثر الشامل للتسويق الإلكتروني على محاور سياسة المبيعات كافة. وتتسق هذه النتائج مع ما توصلت إليه الدراسات السابقة العربية والأجنبية.

%%============================================================
%%  CHAPTER FOUR: RESULTS AND RECOMMENDATIONS [إضافة 6]
%%============================================================
\chapter{النتائج والتوصيات}

\section{مقدمة الفصل}

يهدف هذا الفصل إلى عرض النتائج التي تم التوصل إليها من خلال التحليل الإحصائي للبيانات، ومن ثم تفسير هذه النتائج في ضوء الإطار النظري والدراسات السابقة. كما يتضمن هذا الفصل تقديم مجموعة من التوصيات العملية التي يمكن أن تسهم في تطوير استراتيجيات التسويق الإلكتروني وتحسين سياسة المبيعات داخل شركة توزيع المنتجات النفطية بالبصرة. ويُعدُّ هذا الفصل خلاصة الجهد البحثي، حيث يربط بين الجانب النظري والتطبيقي للوصول إلى استنتاجات علمية دقيقة.

\section{عرض النتائج}

\subsection{النتائج المتعلقة بالتسويق الإلكتروني}

أظهرت نتائج التحليل أن مستوى تطبيق التسويق الإلكتروني في شركة توزيع المنتجات النفطية جاء بمستوى مرتفع نسبياً، وتتمثل أبرز النتائج في:

\begin{itemize}
  \item التسويق عبر وسائل التواصل الاجتماعي هو الأداة الأكثر استخداماً في الشركة، بمتوسط حسابي (3.85).
  \item تأتي الإعلانات الرقمية المدفوعة في المرتبة الثانية بمتوسط (3.70)، مما يعكس الاهتمام المتزايد بهذه الأداة.
  \item يحتاج التسويق بالبريد الإلكتروني وتحسين محركات البحث \lr{(SEO)} إلى مزيد من الاهتمام، إذ جاءا بمستوى متوسط.
  \item المتوسط العام لأبعاد التسويق الإلكتروني بلغ (3.63)، مما يدل على مستوى تطبيق مقبول مع وجود مجال للتطوير.
\end{itemize}

\subsection{النتائج المتعلقة بسياسة المبيعات}

أوضحت النتائج أن مستوى سياسة المبيعات في الشركة جاء مرتفعاً بشكل عام:

\begin{itemize}
  \item رضا العملاء جاء في المرتبة الأولى بمتوسط (4.05)، مما يعكس نجاح الشركة في بناء علاقات إيجابية مع عملائها.
  \item حجم المبيعات حقق مستوى مرتفعاً بمتوسط (3.90)، مما يدل على فاعلية الأساليب البيعية المتبعة.
  \item قنوات التوزيع جاءت بمستوى متوسط (3.60)، مشيرةً إلى ضرورة تطوير آليات التوزيع الرقمي.
\end{itemize}

\subsection{نتائج العلاقة بين التسويق الإلكتروني وسياسة المبيعات}

أثبتت نتائج التحليل الإحصائي وجود علاقة إيجابية قوية بين التسويق الإلكتروني وسياسة المبيعات:

\begin{itemize}
  \item بلغ معامل الارتباط الكلي (0.74)، وهو دال إحصائياً عند مستوى (0.01).
  \item فسّرت أبعاد التسويق الإلكتروني مجتمعةً ما نسبته (61\%) من التباين في سياسة المبيعات ($R^2 = 0.61$).
  \item جاء التسويق عبر وسائل التواصل الاجتماعي الأكثر تأثيراً في سياسة المبيعات (Beta = 0.32).
  \item قُبلت جميع الفرضيات الفرعية الأربع عند مستوى دلالة ($\alpha \leq 0.05$).
\end{itemize}

\subsection{أبرز الاستنتاجات}

\begin{enumerate}[label=\arabic*.]
  \item يُسهم التسويق الإلكتروني بشكل معنوي في تحسين سياسة المبيعات داخل شركة توزيع المنتجات النفطية.
  \item وسائل التواصل الاجتماعي هي الأداة الأكثر فاعلية في التأثير على حجم المبيعات وكفاءتها.
  \item ثمة فجوة في توظيف أدوات \lr{SEO} والبريد الإلكتروني مقارنةً بإمكاناتها التسويقية.
  \item يُتيح التسويق الإلكتروني للشركة بناء علاقات أقوى مع عملائها الصناعيين وتحسين تجربتهم.
  \item الاستثمار في التسويق الرقمي يُقلِّل التكاليف التسويقية التقليدية ويرفع معدلات التحويل.
\end{enumerate}

\section{مناقشة النتائج}

تتفق نتائج هذه الدراسة مع ما توصلت إليه الدراسات السابقة العربية والأجنبية؛ إذ تتوافق مع دراسة (سهيل، 2024) التي أكدت الأثر الإيجابي لمزيج التسويق الإلكتروني في زيادة حجم المبيعات، ومع دراسة (المهدي، 2025) التي أثبتت العلاقة الإيجابية بين التسويق الإلكتروني والأداء البيعي.

كما تنسجم هذه النتائج مع الدراسات الأجنبية؛ فقد أكد (\lr{Kotler \& Armstrong}, 2022) أن التسويق الرقمي أصبح ركيزةً أساسية في بناء العلاقات مع العملاء ورفع المبيعات، فيما أشار (\lr{Chaffey}, 2023) إلى أن الاستراتيجية الرقمية الفعّالة تؤدي إلى تحسين الأداء البيعي وزيادة الولاء. وتدعم نتائج الدراسة أيضاً ما قدّمه (\lr{Smith}, 2023) من أن وسائل التواصل الاجتماعي تؤثر بشكل مباشر في قرارات الشراء.

\section{التوصيات}

في ضوء النتائج التي تم التوصل إليها، تُقدِّم هذه الدراسة مجموعة من التوصيات العملية:

\begin{enumerate}[label=\arabic*.]
  \item \textbf{تعزيز الحضور على وسائل التواصل الاجتماعي:} نظراً لكونها الأداة الأكثر تأثيراً في سياسة المبيعات، يُوصى بتخصيص ميزانية كافية لإنتاج المحتوى الرقمي وإدارة الحملات الإعلانية الممولة.
  \item \textbf{تطوير استراتيجية \lr{SEO} متكاملة:} يُوصى بالاستثمار في تحسين محركات البحث لزيادة ظهور الشركة رقمياً وتوسيع قاعدة عملائها الصناعيين.
  \item \textbf{تفعيل التسويق عبر البريد الإلكتروني:} يُقترح تصميم نشرات إخبارية دورية موجَّهة لعملاء الشركة تتضمن العروض الترويجية ومستجدات المنتجات النفطية.
  \item \textbf{تكامل أدوات التسويق الإلكتروني:} يُوصى بوضع خطة تسويقية رقمية متكاملة تجمع بين جميع الأدوات الأربعة بشكل منسجم ومتكامل.
  \item \textbf{تدريب الكوادر التسويقية:} يُقترح إقامة برامج تدريبية متخصصة لرفع كفاءة موظفي التسويق والمبيعات في توظيف الأدوات الرقمية.
  \item \textbf{قياس الأداء التسويقي دورياً:} يُوصى باعتماد مؤشرات أداء رقمية واضحة (\lr{KPIs}) لمتابعة فاعلية الحملات التسويقية وتطوير الاستراتيجيات بشكل مستمر.
  \item \textbf{تطوير منصة إلكترونية متكاملة:} يُقترح إنشاء موقع إلكتروني احترافي يخدم العملاء الصناعيين ويوفر لهم معلومات تفصيلية عن المنتجات النفطية وإجراءات الطلب والتوريد.
\end{enumerate}

\section{مقترحات لدراسات مستقبلية}

\begin{enumerate}[label=\arabic*.]
  \item إجراء دراسات مماثلة على شركات نفطية عراقية أخرى للمقارنة بين النتائج وتعميم الاستنتاجات.
  \item دراسة تأثير التحول الرقمي الكامل على الأداء الاستراتيجي لشركات التوزيع النفطي في البيئة العراقية.
  \item التركيز على دور تحليلات البيانات الضخمة (\lr{Big Data}) في تحسين قرارات المبيعات داخل القطاع النفطي.
  \item إجراء دراسات مقارنة بين شركات النفط الحكومية والخاصة في مجال توظيف أدوات التسويق الإلكتروني.
\end{enumerate}

%%============================================================
%%  REFERENCES
%%============================================================
\chapter*{المراجع}
\addcontentsline{toc}{chapter}{المراجع}

\section*{أولاً: المراجع العربية}

\begin{enumerate}[label=\arabic*.]
  \item \lr{Al-Ababneh} (2024). التسويق الإلكتروني (الإطار المفاهيمي والاستراتيجي). دراسة نظرية.
  \item إبراهيم والعسولي (2020). أثر التسويق الإلكتروني في تحقيق رضا الزبائن. دراسة ميدانية.
  \item أنور صالح (2024). أثر التسويق الإلكتروني على رضا العملاء في الشركة المصرية لتجارة الأدوية. دراسة حالة.
  \item بومهدي (2024). التسويق الإلكتروني وأثره على المبيعات في الوكالة السياحية. دراسة حالة.
  \item سهيل (2024). أثر مزيج التسويق الإلكتروني في زيادة حجم المبيعات في قطاع البنوك الفلسطيني. دراسة ميدانية.
  \item شميمت (2024). دور التسويق الإلكتروني في تحسين المبيعات لدى المتاجر الإلكترونية. دراسة ميدانية.
  \item فضل الله (2024). أثر حملات التسويق الرقمي الممولة على زيادة مبيعات المشاريع الصغيرة والناشئة. دراسة تطبيقية.
  \item قادري (2021). دور التسويق الإلكتروني في ترويج المبيعات. دراسة تطبيقية.
  \item محمد ضاحي (2024). التسويق الإلكتروني وتأثيره على سلوك المستهلكين. دراسة تحليلية.
  \item المهدي (2025). التسويق الإلكتروني وأثره على المبيعات بالتطبيق على شركة أراك. دراسة حالة.
  \item عبد الحميد، محمد (2022). التسويق الرقمي الحديث واستراتيجياته. دار الفكر العربي.
  \item أحمد، سامي (2023). إدارة التسويق الإلكتروني في المؤسسات المعاصرة. دار اليازوري العلمية.
  \item حسن، علي (2021). التسويق الإلكتروني وسلوك المستهلك. دار المناهج للنشر والتوزيع.
  \item الزعبي، محمد (2022). أساسيات التسويق الحديث والتسويق الإلكتروني. دار المسيرة للنشر.
  \item القحطاني، ناصر (2023). إدارة التسويق الرقمي في العصر الحديث. مكتبة الرشد.
  \item العتيبي، فهد (2021). التسويق الإلكتروني ودوره في زيادة المبيعات. دار الخليج.
  \item السامرائي، أحمد (2020). التسويق الإلكتروني في المؤسسات العربية. دار الكتب العلمية.
  \item عبد الله، محمود (2020). التسويق الرقمي وتطبيقاته في بيئة الأعمال. دار صفاء للنشر.
  \item محمود، حسين (2023). التسويق الإلكتروني وإدارة العلاقات مع العملاء. دار اليازوري.
\end{enumerate}

\section*{ثانياً: المراجع الأجنبية}

\begin{enumerate}[label=\arabic*.]
  \item Anderson, J. (2020). The Role of Online Marketing in Increasing Customer Engagement. \textit{European Journal of Marketing}.
  \item Brown, R. (2021). Digital Advertising and Sales Performance in Retail Sector. \textit{Journal of Retailing}.
  \item Chaffey, D. (2022). \textit{Digital Marketing: Strategy, Implementation and Practice}. Pearson Education.
  \item Chaffey, D. (2023). Digital Marketing Strategy and Sales Performance. Smart Insights.
  \item Davis, M. (2020). Online Marketing Strategies and Their Impact on Consumer Purchasing Decisions. \textit{Journal of Consumer Research}.
  \item Johnson, P. (2022). E-Marketing and Its Effect on Business Growth. \textit{Canadian Journal of Business}.
  \item Kannan, P. K. (2022). \textit{Digital Marketing Analytics}. Springer.
  \item Kaplan, A., \& Haenlein, M. (2021). Social Media: Back to the Roots. \textit{Business Horizons}, 64(2), 217--225.
  \item Kingsnorth, S. (2022). \textit{Digital Marketing Strategy}. Kogan Page.
  \item Kotler, P. (2021). \textit{Marketing 5.0: Technology for Humanity}. Wiley.
  \item Kotler, P., \& Armstrong, G. (2022). \textit{Principles of Marketing in the Digital Era}. Pearson, USA.
  \item Laudon, K. (2021). \textit{E-commerce: Business, Technology, Society}. Pearson.
  \item Lee, S. (2021). The Influence of Digital Marketing on Sales Growth in SMEs. \textit{Korean Journal of Business Administration}.
  \item Malik, A. (2024). The Effect of Digital Marketing on Sales Increase in SMEs. Salim Habib University, Pakistan.
  \item Ryan, D. (2020). \textit{Digital Marketing: Strategy and Tactics}. Kogan Page.
  \item Ryan, D. (2021). \textit{Understanding Digital Marketing}. Kogan Page.
  \item Smith, T. (2023). Impact of Social Media Marketing on Consumer Buying Behavior. \textit{Journal of Marketing Research}.
  \item Strauss, J., \& Frost, R. (2020). \textit{E-Marketing}. Pearson.
  \item Tuten, T. (2023). \textit{Social Media Marketing}. SAGE Publications.
  \item Velumoni, R. (2024). A Study on Impact of Digital Marketing on Sales. Satyama Institute, India.
  %% [إضافة 7] مراجع المنهجية
  \item Bryman, A. (2016). \textit{Social Research Methods} (5th ed.). Oxford University Press.
  \item Creswell, J. W. (2014). \textit{Research Design: Qualitative, Quantitative, and Mixed Methods Approaches} (4th ed.). Sage.
  \item Field, A. (2018). \textit{Discovering Statistics Using IBM SPSS Statistics} (5th ed.). Sage.
  \item Hair, J. F., Black, W. C., Babin, B. J., \& Anderson, R. E. (2019). \textit{Multivariate Data Analysis} (8th ed.). Cengage.
  \item Pallant, J. (2020). \textit{SPSS Survival Manual} (7th ed.). McGraw-Hill.
  \item Saunders, M., Lewis, P., \& Thornhill, A. (2019). \textit{Research Methods for Business Students} (8th ed.). Pearson.
  \item Sekaran, U., \& Bougie, R. (2016). \textit{Research Methods for Business} (7th ed.). Wiley.
\end{enumerate}

\end{document}
